\documentclass[11pt]{ctexart}
\usepackage[a4paper,margin=1in]{geometry}
\usepackage{graphicx}
\usepackage{xcolor}
\usepackage{hyperref}
\definecolor{refgray}{gray}{0.45}
\title{\textbf{Market Views｜全球投行叙事汇编}\\[0.4em]{\large 通胀传导分化与AI资本开支加速主导全球市场，央行观望与政策分歧加剧不确定性}}
\author{KC桌面 自动生成}
\date{260726}
\begin{document}
\maketitle
\tableofcontents
\newpage
\section{一页摘要}
\begin{itemize}
  \item 全球通胀压力从上游向下游传导力度弱于2022年，但供应链瓶颈和服务业价格粘性仍构成不确定性；主要央行普遍处于观望状态，美联储7月料维持利率不变，内部分歧聚焦于通胀回落速度是否足够快。
  \item AI基础设施投资正从‘追赶需求’转向‘主动超前配置’，Google连续上调资本开支指引至1950-2050亿美元，带动半导体设备、数据中心及测试环节需求结构性增长；AI应用软件收入拐点初现，ServiceNow AI相关年化合同价值突破10亿美元。
  \item 具身智能已跨过量产拐点，但需求结构仍以非生产性用途为主，高质量物理交互数据成为制约‘大脑’成熟的核心瓶颈；AI芯片测试与电源架构升级为测试设备商带来结构性增长机遇。
  \item 中国生物科技板块进入业绩分化期：康方生物与科伦博泰凭借核心产品放量及一次性收益实现盈利转正，而药明系受汇率与并购费用拖累利润增速；辉瑞估值处于历史低位，但市场焦点已转向下半年关键临床数据。
  \item 消费与汽车板块呈现显著结构性分化：啤酒行业总量趋平但企业间份额加速转移，华润啤酒表现突出；汽车零部件领域，潍柴动力的估值逻辑正与AI硬件主题挂钩，恒立液压的滚珠丝杠订单远超指引。
  \item 全球工业与机器人领域呈现结构性分化：机器人手术渗透率加速提升但竞争格局松动，电梯行业订单稳健但成本通胀侵蚀利润，中国钢铁供需背离、库存被动累积。
  \item 欧洲电动化进程在2026年上半年显著加速，BEV注册量同比增35\%，但充电基础设施季度建设速度同比下降20\%，供需缺口扩大；中国电池龙头宁德时代以创纪录业绩和回购计划强化市场信心。
  \item 市场对香港银行净息差走势的预期出现关键修正，从‘逐年收窄’转向‘持稳’，叠加信用成本下降和股东回报增强预期，中银香港等标的的盈利预测和估值逻辑面临重估。
\end{itemize}
\section{全球投行叙事汇编}
今日纳入 59 篇投行/券商研究，按 10 个市场、资产与产业主题系统整理。
\newpage
\subsection{宏观与利率：通胀传导分化与政策观望}
\textbf{全球通胀压力从上游向下游的传导力度弱于2022年，但供应链瓶颈和服务业价格粘性仍构成不确定性；主要央行普遍处于观望状态，美联储7月料维持利率不变，但内部分歧聚焦于通胀回落速度是否足够快；中国宏观活动分化，人民币升值节奏放缓但趋势未改。}
\subsubsection*{主流共识}
\begin{itemize}
  \item 主要央行普遍处于观望状态，美联储7月FOMC料维持利率不变，欧洲央行本周会议维持利率不变，9月加息已被定价。
  \item 全球PPI向CPI的传导力度弱于2022年，企业利润缓冲更充足、消费需求更温和、财政货币支持减弱是主因。
  \item 中国宏观活动呈现分化：生产端指标保持韧性，但消费信心和部分工业需求出现边际走弱。
\end{itemize}
\subsubsection*{分歧与条件差异}
\begin{itemize}
  \item 美联储内部对‘当前利率是否足够限制性’尚未统一：副主席Jefferson和纽约联储Williams认为政策定位良好，而达拉斯联储Logan和克利夫兰联储Hammack认为可能需要更高利率。
  \item 通胀粘性的归因存在分歧：部分官员强调关税、能源和AI投资等结构性因素，而另一些官员认为临时性因素放大了6月数据表现。
  \item 欧元区面临全球能源冲击和利差分化双重压力，欧元失去独立定价能力；而日本成为发达市场中少见的资金流入目的地，可能对日元形成支撑。
\end{itemize}
\subsubsection*{机构观点}
\paragraph{BARC} 人民币升值节奏放缓并非趋势逆转，政策意图是控制变化节奏而非制造弱势货币。央行持续通过中间价管理变化节奏，中间价持续弱于市场预期；CFETS人民币汇率指数接近2022年8月以来高位，贸易竞争力存在底线；出口商结汇意愿从高位回落但仍处历史高位。维持2026年底美元兑人民币6.65预测，升值路径将更为渐进。
{\small\color{refgray}数据：2026年底美元兑人民币预测6.65\par}
{\small\color{refgray}数据：CFETS人民币汇率指数接近2022年8月以来最高水平\par}
{\small\color{refgray}数据：3个月风险逆转指标的人民币看多程度较年初减弱\par}
{\small\color{refgray}边际变化：近期人民币在突破6.80后进入盘整，升值节奏放缓，但方向未变；中间价与市场隐含水平的差距收窄但未消失。\par}
\paragraph{NOM} 2026年PPI向CPI的传导力度明显弱于2022年，且各国之间差异显著。全球CPI-PPI差值在3月转负后4-5月进一步扩大，显示企业利润空间被压缩，但初始缓冲强于2022年。若能源价格持续上行，企业可能逐步调整定价策略。高暴露地区包括新加坡、秘鲁、智利、中国台湾、加拿大。
{\small\color{refgray}数据：2026年PPI与CPI累计涨幅相关性远低于2022年\par}
{\small\color{refgray}数据：全球CPI-PPI差值2026年3月转负，4-5月进一步扩大\par}
{\small\color{refgray}数据：27国中13个高暴露地区被识别\par}
{\small\color{refgray}边际变化：本轮能源价格上涨中企业利润缓冲更充足，无需立即激进转嫁成本；但利润空间正在被压缩，若能源价格持续上行，企业可能被迫调整定价。\par}
\paragraph{NOM} 美联储7月FOMC将维持政策利率不变，但可能出现两名成员投出加息异议票。6月核心PCE通胀预计回落至月环比0.175\%，折年率接近2\%，与美联储2\%通胀目标取得实质性进展。市场对7月加息的定价更多反映前瞻指引不足和反应函数不确定性，而非宏观前景实质性变化。
{\small\color{refgray}数据：6月核心PCE月环比预计0.175\%，折年率约2\%\par}
{\small\color{refgray}数据：6月核心CPI月环比下降0.017\%，为2020年5月以来首次负增长\par}
{\small\color{refgray}数据：2026年底核心PCE同比预计降至3.2\%（考虑方法修订约3.0\%）\par}
{\small\color{refgray}边际变化：通胀降温信号较为明确，但6月数据可能受临时因素放大；多数官员转向‘需要看到多少证据才能确认通胀回落’。\par}
\paragraph{GS} 中国7月宏观活动呈现分化：30城新房日均成交面积回升至接近去年同期水平，16城二手房日均成交面积继续高于去年同期；但Morning Consult消费者信心指数下滑，交通拥堵指数回落，能源成本上升可能对消费信心产生影响。生产端钢铁产量增加但需求边际下降，沿海省份煤炭日耗环比回落但仍高于去年同期。
{\small\color{refgray}数据：30城新房日均成交面积接近去年同期水平\par}
{\small\color{refgray}数据：16城二手房日均成交面积继续高于去年同期\par}
{\small\color{refgray}数据：Morning Consult消费者信心指数过去一周下滑\par}
{\small\color{refgray}数据：汽油柴油价格7月17日分别上调300元/吨和290元/吨\par}
{\small\color{refgray}边际变化：消费信心与拥堵指数同步走弱，叠加能源价格上涨，提示7月下旬消费动能面临短期不确定性。\par}
\paragraph{GS} 美联储决策框架重心正从‘利率水平是否足够限制性’转向‘通胀回落速度是否足够快’。6月CPI数据后，多数官员不再争论是否加息，而是转向需要多少证据确认通胀回落。但内部分歧未消：副主席Jefferson和纽约联储Williams认为政策定位良好，达拉斯联储Logan和克利夫兰联储Hammack认为可能需要更高利率。通胀粘性归因从能源转向关税、AI投资等结构性因素。
{\small\color{refgray}数据：6月核心CPI月环比下降0.017\%\par}
{\small\color{refgray}数据：达拉斯联储Logan表示‘适度提高利率能更好平衡双重使命不确定性’\par}
{\small\color{refgray}数据：克利夫兰联储Hammack认为通胀持续偏高可能需要更高利率\par}
{\small\color{refgray}边际变化：官员表态从‘要不要加息’转向‘需要看到多少证据才能确认通胀回落’；通胀粘性归因从能源故事转向结构性因素。\par}
\paragraph{GS} 美国财政部半年度外汇报告未将任何主要贸易伙伴列为汇率操纵国，监测名单维持不变。对中国措辞有所缓和：不再强调‘及时有序推动人民币升值’，转而承认‘当局允许人民币渐进、有管理地升值’；删除‘抑制进口’‘过度依赖出口增长’等表述。泰国、新加坡、瑞士仅满足一项标准，可能被移出监测名单。
{\small\color{refgray}数据：监测名单包含10个地区：中国、日本、韩国、台湾、新加坡、越南、德国、爱尔兰、泰国、瑞士\par}
{\small\color{refgray}数据：泰国、新加坡、瑞士仅满足三项标准中的一项\par}
{\small\color{refgray}数据：新关税从Section 122转向Section 301，有效税率仅上升1个百分点至8\%\par}
{\small\color{refgray}边际变化：对中国措辞有所缓和，不再强调‘及时有序推动人民币升值’；美元走弱使报告整体基调更温和。\par}
\paragraph{GS} 外汇市场表面平静，驱动回报表现的仍是利差交易和能源贸易条件。欧元与美元相关系数异常紧密，失去独立定价能力，受全球能源冲击和利差分化双重压力。英镑面临财政审查回归与加息定价修正双重不确定性。日元空头仓位接近7月干预前水平，但当前美国经济距离衰退较远，若AI估值受质疑可能成为日元大幅走强催化剂。
{\small\color{refgray}数据：欧元与广义美元相关系数超过其他任何G10货币对\par}
{\small\color{refgray}数据：USD/JPY升至40年新高\par}
{\small\color{refgray}数据：日本权益类产品四周累计净流入约65亿美元\par}
{\small\color{refgray}边际变化：能源价格高企为美元提供缓冲；日本国内资金从海外回流本土资产，可能对日元形成支撑。\par}
\paragraph{GS} 7月全球制造业PMI显示整体活动扩张，但供应链端不确定性未消退。发达市场综合PMI从51.2升至52.8，制造业供应商交付时间指数虽略有缓解，仍处于历史高位附近（41.8），远低于50荣枯线。制造业与服务业价格走势分化：制造业投入价格PMI从71.7降至68.4，服务业投入价格PMI从62.0升至62.4，通胀从商品端向服务端传导。
{\small\color{refgray}数据：发达市场综合PMI从51.2升至52.8\par}
{\small\color{refgray}数据：制造业供应商交付时间指数41.8，远低于50荣枯线\par}
{\small\color{refgray}数据：美国交付时间指数42.0，欧元区40.5，日本39.2\par}
{\small\color{refgray}数据：制造业投入价格PMI从71.7降至68.4\par}
{\small\color{refgray}边际变化：交付时间指数仍处于历史高位附近，供应链压力未完全缓解；制造业价格压力缓解，服务业价格压力上升。\par}
\paragraph{GS} 截至7月22日当周，全球权益类产品净流入约300亿美元，较前一周560亿美元收窄。日本成为发达市场少见的资金流入目的地（净流入约15亿美元），美国净流出约72亿美元。新兴市场由中国内地、韩国和台湾带动净流入约296亿美元。科技板块仍为权益类最大流入领域，但流入速度从高位回落。日本债券产品近期出现国内资金流入上升趋势。
{\small\color{refgray}数据：全球权益类产品净流入约300亿美元（前一周560亿美元）\par}
{\small\color{refgray}数据：日本权益类产品净流入约15亿美元，四周累计约65亿美元\par}
{\small\color{refgray}数据：美国权益类产品净流出约72亿美元\par}
{\small\color{refgray}数据：新兴市场权益类产品净流入约296亿美元\par}
{\small\color{refgray}边际变化：日本成为发达市场少见的资金流入目的地；日本债券产品国内资金流入上升趋势早于官方鼓励国内投资表态。\par}
\paragraph{DB} 通力2026年第二季度订单量同比增长11\%，现代化改造业务增速超过15\%，但市场反应被订单利润率小幅下滑和集团利润率扩张有限所抑制。这份DB研报的核心判断是：通力在执行层面持续表现良好，其市场领导地位反映在强劲的订单增长中，但行业正面临成本通胀上升的挑战，公司正通过提价来维持价格成本稳定。
{\small\color{refgray}数据：Q2订单增长中，除中国外所有区域均实现提价，中国定价保持稳定。管理层指出，从招标到订单的确认周期存在明显滞后，这导致订单利润率出现小幅下降。维护业务方面，随着中国组合观察管理的影响逐渐消退，预计收入增速将略有提升。现代化改造业务订单增长强劲，但收入增速仅为7\%，低于正常水平。报告预计，下半年现代化改造收入将重新加速，因为同期基数较低。\par}
{\small\color{refgray}数据：KC评论： 订单增长与收入确认之间的时间差，是理解通力短期财务表现的关键。Q2订单利润率下滑并非需求出现变化，而是招标到订单的确认节奏问题，这一点在评估公司基本面时容易被忽略。\par}
\subsubsection*{关键数据}
\begin{itemize}
  \item 6月核心PCE月环比预计0.175\%，折年率约2\%
  \item CFETS人民币汇率指数接近2022年8月以来最高水平
  \item 全球CPI-PPI差值2026年3月转负，4-5月进一步扩大
  \item 发达市场制造业供应商交付时间指数41.8，远低于50荣枯线
  \item 日本权益类产品四周累计净流入约65亿美元
  \item 30城新房日均成交面积接近去年同期水平
  \item Morning Consult消费者信心指数过去一周下滑
  \item 欧元与广义美元相关系数超过其他任何G10货币对
\end{itemize}
\begin{figure}[htbp]
\centering
\includegraphics[width=0.88\linewidth]{figures/fig_011.jpg}
\caption{Figure 1 - Audrey Ong BARC Bank, Singapore FIGURE 1. CNY has been fixed consistently weaker than the market over recent months FIGURE 2. CFETS CNY has appreciated over recent months 2) The CNY's relative strength is becoming a consideration. The CFETS CNY trade-weighted }
\end{figure}
\begin{figure}[htbp]
\centering
\includegraphics[width=0.88\linewidth]{figures/fig_043.jpg}
\caption{Figure 1 - \textbackslash{}- Initial conditions of healthier margins, weaker consumer demand and less fiscal and monetary support help explain the less aggressive cost pass-through this year. PPI inflation measures upstream prices in the supply chain and is generally a leading indicato}
\end{figure}
\begin{figure}[htbp]
\centering
\includegraphics[width=0.88\linewidth]{figures/fig_080.jpg}
\caption{Exhibit 1 - Exhibit 1: Manufacturing Suppliers' Delivery Times Eased in ex-US DMs but Remain Near All-Time Highs \#\# Elevated Manufacturing Suppliers' Delivery Times \#\# DM Trends Exhibit 2: The DM Composite Flash PMI Rose by +1.6pt to 52.8 in}
\end{figure}
\begin{figure}[htbp]
\centering
\includegraphics[width=0.88\linewidth]{figures/fig_078.jpg}
\caption{Exhibit 1 - Exhibit 1: The Euro's correlation to the broad Dollar has been tighter than for any other G10 Dollar pair GBP: Week one weakness. Spending announcements in the first week of the Burnham government have been relatively minor in sc}
\end{figure}
{\small\color{refgray}本节综合 10 篇研究；涉及 BARC、NOM、DB、GS。}
\newpage
\subsection{AI基础设施与应用：资本开支加速，测试与软件变现进入验证期}
\textbf{AI基础设施投资正从‘追赶需求’转向‘主动超前配置’，Google连续上调资本开支指引至1950-2050亿美元，带动半导体设备、数据中心及测试环节需求结构性增长。同时，AI在应用软件（ServiceNow）和药物研发（AIDD）领域已出现可量化的收入拐点，但传统安全工具面对LLM驱动的攻击暴露短板，催生新的防御需求。}
\subsubsection*{主流共识}
\begin{itemize}
  \item AI基础设施资本开支持续加速，Google上调全年指引至1950-2050亿美元，下半年环比增长约50\%，服务器支出占比升至60\%。
  \item 半导体后端测试设备受益于AI芯片复杂度提升，测试工序增加（老化测试、系统级测试）推高单芯片测试价值，探针卡和同轴插座业务预计年复合增长率达67\%-86\%。
  \item AI应用软件收入拐点初现，ServiceNow AI相关年化合同价值突破10亿美元且加速增长，代理部署量九个月增长九倍。
  \item 数据中心需求强劲，Digital Realty季度签约中约20\%与AI部署直接相关，开发管线超200亿美元，63\%容量已预租。
\end{itemize}
\subsubsection*{分歧与条件差异}
\begin{itemize}
  \item 硅光子（SiPh）设备需求规模：Bernstein认为Disco的硅光子业务正在向量产过渡，收入来源多元但规模仍小；JPM则强调其仍处于研发阶段，对短期出货贡献有限。
  \item 中国OSAT投资节奏：Bernstein指出Q1 OSAT占比超预期后Q2预计回落，因中国OSAT研发活动调整；NOM则认为AI OSAT资本支出增加是结构性增长动力。
  \item 传统安全工具有效性：MS认为传统SIEM/XDR面对LLM驱动的攻击存在明显缺口，需结合行为基线和XDR+SOAR；但部分企业认为现有工具通过升级仍可应对。
  \item SAP利润率短期波动：Bernstein认为利润率低于预期是研发投入和收购稀释所致，短期可控；市场则担忧利润率趋势能否持续改善。
\end{itemize}
\subsubsection*{机构观点}
\paragraph{Bernstein} Disco Q2出货指引1410亿日元超预期，HBM需求强劲，硅光子成为新增长来源，但规模仍小于HBM/CoWoS；Q1毛利率71.3\%超指引，耗材利润率改善趋势延续。SAP 2Q26 CCB增速26\%超预期，缓解市场对增长节奏的担忧，运营利润率27.8\%低于预期但属短期稀释。AMD Helios架构发布，Anthropic成为第三大客户，承诺研究50亿美元，CPU与加速器2030年TAM大幅上调至2万亿美元。
{\small\color{refgray}数据：Disco Q2出货指引1410亿日元，环比+4\%，高于市场预期的1311亿日元\par}
{\small\color{refgray}数据：SAP CCB达229.3亿欧元，固定汇率同比+26\%，高于云收入增速24\%\par}
{\small\color{refgray}数据：AMD将2030年CPU TAM从1200亿上调至2200亿美元，加速器TAM看至1.4万亿美元\par}
{\small\color{refgray}数据：Anthropic计划部署2GW算力，首个1GW于2027上半年开始部署\par}
{\small\color{refgray}边际变化：硅光子从研发向量产过渡，SAP CCB增速稳定缓解市场核心担忧，AMD从GPU追赶者转向CPU+GPU双增长曲线平台供应商。\par}
\paragraph{NOM} 鸿劲精密（Hon. Precision）受益于AI/HPC芯片测试需求结构性增长，2026-2028年营收复合增长率预计达59\%，ATC升级路径清晰，单机价值量存在上行空间。另一家测试接口厂商探针卡和同轴插座业务预计分别实现67\%和86\%的年复合增长率，系统级测试与老化测试成为AI芯片新增测试工序，共封装光学测试提供额外增长空间。
{\small\color{refgray}数据：鸿劲精密2026-2028年营收CAGR预计59\%，EPS从134元增长至336元\par}
{\small\color{refgray}数据：探针卡业务2026-2028年收入CAGR 67\%，同轴插座CAGR 86\%\par}
{\small\color{refgray}数据：AI/HPC芯片占测试设备订单约80\%\par}
{\small\color{refgray}数据：下一代GPU平台采用GPU-on-GPU SoIC堆叠，推动ATC升级\par}
{\small\color{refgray}边际变化：AI芯片测试工序从传统探针卡/插座延伸至老化测试、系统级测试和CPO测试，测试价值量结构性提升。\par}
\paragraph{Citi} Digital Realty二季度现金续约价差达25\%，大于1MW批发租赁续约价差高达67\%；季度签约1.08亿美元中约20\%与AI部署直接相关，AI相关季度签约量有望维持1亿美元或更高。开发管线超200亿美元，63\%容量已预租，积压订单创纪录达19亿美元，相当于运营收入的30\%。
{\small\color{refgray}数据：现金续约价差25\%，大于1MW租赁续约价差67\%\par}
{\small\color{refgray}数据：二季度签约1.08亿美元，约20\%直接来自AI部署\par}
{\small\color{refgray}数据：开发管线超200亿美元，预期回报率11.5\%，63\%已预租\par}
{\small\color{refgray}数据：积压订单19亿美元，占运营收入30\%\par}
{\small\color{refgray}边际变化：AI需求从概念转化为具体订单，管理层将核心运营指标增长预期从高个位数上调至双位数。\par}
\paragraph{JPM} Google将2026年全年资本开支指引从1800-1900亿美元上调至1950-2050亿美元，连续第二次上调，服务器支出占比升至60\%。Google Cloud收入同比增长82\%，订单积压突破5000亿美元，未来两年内50\%将确认为收入。Disco Q1营业利润490亿日元符合预期，先进封装研发工具约30亿日元延迟确认，总部生产支援团队继续保留反映需求稳健。
{\small\color{refgray}数据：Google 2026年资本开支指引上调至1950-2050亿美元，下半年环比增长约50\%\par}
{\small\color{refgray}数据：Google Cloud订单积压5140亿美元，环比增加超500亿美元\par}
{\small\color{refgray}数据：Disco Q1营业利润490亿日元，出货额1359亿日元，环比+12\%\par}
{\small\color{refgray}数据：Disco约30亿日元先进封装研发工具延迟至后续季度确认\par}
{\small\color{refgray}边际变化：Google资本开支从‘追赶需求’转向‘主动超前配置’，且预计明年仍将显著增长；Disco生产支援团队保留暗示产能瓶颈缓解慢于预期。\par}
\paragraph{MS} ServiceNow 2Q26 cRPO同比增长21.5\%超预期，订阅收入同比增长23\%，AI相关年化合同价值突破10亿美元且加速增长，代理部署量九个月增长九倍。Hugging Face遭LLM自主代理攻击，传统SIEM依赖静态规则难以检测，需结合行为基线和XDR+SOAR才能有效防御，组合工具可将平均修复时间缩短96\%。
{\small\color{refgray}数据：ServiceNow cRPO固定汇率同比+21.5\%，超预期约2个百分点\par}
{\small\color{refgray}数据：AI相关ACV突破10亿美元，季度内呈加速趋势\par}
{\small\color{refgray}数据：代理部署量九个月增长九倍\par}
{\small\color{refgray}数据：Hugging Face攻击产生超17000条事件记录，最终由LLM分析代理重构\par}
{\small\color{refgray}边际变化：ServiceNow成为应用软件最早看到AI收入拐点的公司之一；传统安全工具面对LLM驱动攻击暴露短板，催生新一代检测需求。\par}
\paragraph{GS} KOSPI上周下跌1.9\%，但外资连续第二周净买入，散户去杠杆但未至恐慌：保证金贷款余额从峰值250亿降至220亿美元，杠杆ETF AUM从530亿腰斩至260亿美元。半导体板块外资持股比例处于近两年标准差以下低位，化工板块盈利上修较强，汽车板块下调最多。
{\small\color{refgray}数据：韩国散户保证金贷款余额从250亿美元降至220亿美元\par}
{\small\color{refgray}数据：杠杆ETF AUM从530亿美元峰值降至260亿美元\par}
{\small\color{refgray}数据：KOSPI半导体板块外资持股比例低于历史均值2个标准差\par}
{\small\color{refgray}数据：KOSPI 12个月远期EPS上周被下调0.4\%\par}
{\small\color{refgray}边际变化：外资连续净买入但半导体板块配置仍轻，散户去杠杆主动减仓而非被动强平，市场情绪从过热回到中性偏低。\par}
\paragraph{JPM} 中国AI驱动药物研发（AIDD）行业正在从业务探索转向早期商业化。JPM在一份覆盖两家相关市场AIDD公司的报告中指出，这一转变有三大支撑：AIDD已实质性地改善了药物研发的生产效率，并通过“干湿实验室”闭环平台进一步提高成功率；中国在研发成本、速度和监管方面的结构性优势，使AIDD应用环境比海外更具吸引力；大型药企的交易流、临床里程碑和收入增长为AIDD公...
{\small\color{refgray}数据：报告覆盖的两家公司——英矽智能（3696.HK）和晶泰科技（2228.HK）——被更新公司观点。报告认为，这两家公司提供了互补性的研究主题：英矽智能是高不确定性、高回报的管线型公司，晶泰科技则是以AIDD服务和机器人湿实验室为核心的多元化业务平台。\par}
{\small\color{refgray}数据：报告提供了两组可核验的效率数据。在晶泰科技的案例中，传统需要约2.5年的先导化合物发现时间被压缩至不到2个月；400分子库的合成时间从至少3周缩短至5天。英矽智能方面，从靶点发现到临床前候选化合物的时间从4.5年降至12至18个月，平均每个PCC的成本控制在300万至500万美元。\par}
\subsubsection*{关键数据}
\begin{itemize}
  \item Google 2026年资本开支指引上调至1950-2050亿美元，下半年环比增长约50\%
  \item ServiceNow AI相关ACV突破10亿美元，代理部署量九个月增长九倍
  \item Digital Realty季度签约中约20\%与AI部署直接相关，积压订单19亿美元
  \item Disco Q2出货指引1410亿日元超预期，硅光子成为新增长来源
  \item 探针卡和同轴插座业务预计2026-2028年CAGR分别达67\%和86\%
  \item AMD将2030年CPU TAM从1200亿上调至2200亿美元
  \item SAP CCB固定汇率同比+26\%，高于云收入增速24\%
  \item 韩国散户杠杆ETF AUM从530亿美元峰值腰斩至260亿美元
\end{itemize}
\begin{figure}[htbp]
\centering
\includegraphics[width=0.88\linewidth]{figures/fig_001.jpg}
\caption{EXHIBIT 4 - EXHIBIT 5: DISCO Q1 revenue was ¥114bn, +27\% yoy EXHIBIT 6: DISCO Q1 shipment was ¥136bn, +22\% yoy EXHIBIT 7: DISCO's Q1 GP / OP margin was 71.3\% / 42.9\%, EXHIBIT 8: DISCO's Q1 EPS was ¥314.57, +44\% yoy +0.4\% / -1.3\% qoq}
\end{figure}
\begin{figure}[htbp]
\centering
\includegraphics[width=0.88\linewidth]{figures/fig_120.jpg}
\caption{Figure 1 - Figure 1: Google's reported capex and qoq trend Figure 2: Google's reported capex and yoy trend Figure 3: Google's free cash flow trend}
\end{figure}
\begin{figure}[htbp]
\centering
\includegraphics[width=0.88\linewidth]{figures/fig_133.jpg}
\caption{Exhibit 1 - Exhibit 1: Q2 cc cRPO Growth of +21.5\% YoY Came in Ahead of Consensus. Organic Growth Remains The Key Focus \#\# What We Liked in Q2: \textbackslash{}- Beat magnitude improved QoQ. NOW returned a more meaningful beat in the quarter on cc subs rev}
\end{figure}
\begin{figure}[htbp]
\centering
\includegraphics[width=0.88\linewidth]{figures/fig_136.jpg}
\caption{Exhibit 1 - Exhibit 1: Traditional SIEM Alone Not Enough For Newer Sophisticated Attacks, but Combination of Tools Leads To Significantly Higher MTTR Exhibit 2: Modern SOC Architecture Combines Detection Engine with Response Exhibit 3: Cus}
\end{figure}
{\small\color{refgray}本节综合 12 篇研究；涉及 Bernstein、NOM、Citi、GS、JPM、MS。}
\newpage
\subsection{科技与AI（2）：具身智能量产拐点与AI测试设备的结构性机遇}
\textbf{具身智能已跨过量产拐点，但需求结构仍以非生产性用途为主，高质量物理交互数据成为制约“大脑”成熟的核心瓶颈；同时，AI芯片测试与电源架构升级为测试设备商带来结构性增长机遇。}
\subsubsection*{主流共识}
\begin{itemize}
  \item 具身智能进入量产阶段，叙事从TAM转向出货量、ASP和实际部署。
  \item AI/HPC芯片测试复杂度提升，SLT设备受益于平台迭代和散热要求升级。
  \item AI服务器电源架构向高压直流（HVDC）迁移趋势明确。
\end{itemize}
\subsubsection*{分歧与条件差异}
\begin{itemize}
  \item 人形机器人出货量验证的是可制造性还是真实需求？NOM认为当前出货量主要反映供给侧能力，而非需求可持续性。
  \item HVDC采用速度存在市场分歧，但NOM认为迁移路径保持完整。
  \item 灵巧手技术路线（连杆、腱绳、直驱）尚未完全收敛，不同方案在成本、寿命和算法友好度上各有优劣。
\end{itemize}
\subsubsection*{机构观点}
\paragraph{NOM} 具身智能已跨过量产拐点，但需求结构仍以非生产性用途为主（娱乐/表演30\%、消费30\%、政府采购20\%、教育15\%），工业/商业仅占5\%。高质量物理交互数据缺口达20倍（2026年需约1000万小时，全球优质数据存量仅50万小时），机器人大脑成熟仍需2-3年。灵巧手三条技术路线正在收敛，连杆手（成本1.2-1.3万）占80\%出货量，腱绳手（约3.8万）寿命仅2个月，直驱手（3.3-3.4万）算法友好但发热问题待解。
{\small\color{refgray}数据：智元机器人累计产量达15000台\par}
{\small\color{refgray}数据：2026年预计出货4.5-5万台，低于媒体所说的10万\par}
{\small\color{refgray}数据：人形机器人日租金同比下降50-60\%\par}
{\small\color{refgray}数据：2026年需约1000万小时高质量交互数据，全球优质数据存量仅50万小时\par}
{\small\color{refgray}边际变化：WAIC 2026机器人展区从杂耍式演示转向实际服务，近60台人形机器人执行真实任务，量产能力得到验证，但需求结构仍以非生产性用途为主，数据瓶颈成为新焦点。\par}
\paragraph{NOM} Chroma同时受益于AI/HPC芯片测试复杂度提升和AI服务器电源架构向HVDC迁移两条结构性主线。SLT设备将随英伟达GPU平台迭代进入新一轮升级周期，2028年GPU-on-GPU SoIC封装方案对热控制提出更严格要求，直接推升测试设备ASP。半导体/光子学测试业务2025-2028年营收CAGR可达68\%。HVDC方面，一家美国头部CSP计划从2026年下半年开始采用±400V直流供电，先应用于VR200机架。
{\small\color{refgray}数据：Chroma未来三年净利润CAGR预计39\%\par}
{\small\color{refgray}数据：半导体/光子学测试业务2025-2028年营收CAGR 68\%\par}
{\small\color{refgray}数据：英伟达计划2028年采用GPU-on-GPU SoIC封装\par}
{\small\color{refgray}数据：一家美国头部CSP计划2026年下半年导入±400V HVDC\par}
{\small\color{refgray}边际变化：市场对HVDC采用速度存在分歧，但NOM认为迁移路径保持完整；SLT设备从可选支出变为必需支出，ASP随封装复杂度提升而上升。\par}
\subsubsection*{关键数据}
\begin{itemize}
  \item 智元机器人累计产量达15000台
  \item 2026年预计出货4.5-5万台，低于媒体所说的10万
  \item 人形机器人日租金同比下降50-60\%
  \item 2026年需约1000万小时高质量交互数据，全球优质数据存量仅50万小时
  \item 连杆手成本1.2-1.3万，占80\%出货量
  \item 腱绳手约3.8万，寿命仅2个月
  \item 直驱手3.3-3.4万，20+自由度
  \item Chroma未来三年净利润CAGR预计39\%
  \item 半导体/光子学测试业务2025-2028年营收CAGR 68\%
  \item 英伟达计划2028年采用GPU-on-GPU SoIC封装
\end{itemize}
{\small\color{refgray}本节综合 2 篇研究；涉及 NOM。}
\newpage
\subsection{中国生物科技：盈利拐点与海外催化剂并行}
\textbf{中国生物科技板块进入业绩分化期：康方生物与科伦博泰凭借核心产品放量及一次性收益实现盈利转正，而药明系受汇率与并购费用拖累利润增速；辉瑞估值处于历史低位，但市场焦点已转向下半年关键临床数据。}
\subsubsection*{主流共识}
\begin{itemize}
  \item 中国生物科技公司核心产品（sac-TMT、卡度尼利单抗、依沃西单抗）进入医保后销售放量，是收入增长的主要驱动力。
  \item 人民币对美元升值（约4\%）对以美元计价收入为主的药明系公司（药明合联、药明生物）利润端形成普遍压力。
  \item 对外授权收入（如康方生物、恒瑞医药）是短期盈利转正或利润弹性的重要来源，但确认节奏存在不确定性。
\end{itemize}
\subsubsection*{分歧与条件差异}
\begin{itemize}
  \item 科伦博泰与康方生物均实现盈利转正，但前者依赖一次性诉讼和解收入（7.5亿元），后者依靠经营利润改善（营业利润率9.7\%），市场对盈利可持续性的评估存在分歧。
  \item 辉瑞估值（2027年PE 8.75倍）处于历史低位，但市场对其下半年关键临床数据（mevrometostat 3期、GLP-1/Amylin早期数据）的成败影响估值幅度（±3\%）存在不同预期。
  \item 恒瑞医药利润创新高（2Q26净利润29亿元），但收入增速仅1\%，利润弹性来自费用压缩与海外资产重估，市场对利润质量与可持续性的看法存在差异。
\end{itemize}
\subsubsection*{机构观点}
\paragraph{NOM} 科伦博泰1H26营收预计14.1亿元（同比+48\%），sac-TMT进入医保后销售加速，毛利率从69.4\%提升至78.5\%；诉讼和解带来7.5亿元其他收入，推动净利润转正至7.31亿元（去年同期亏损1.45亿元）。剔除一次性收益后经营利润约1.13亿元，经营利润率8\%。下半年关注默沙东BLA申报进展。
{\small\color{refgray}数据：1H26营收14.1亿元，同比+48\%\par}
{\small\color{refgray}数据：sac-TMT销售6.5亿元，合作收入7.5亿元\par}
{\small\color{refgray}数据：诉讼和解收入7.5亿元\par}
{\small\color{refgray}数据：净利润7.31亿元（去年同期亏损1.45亿元）\par}
{\small\color{refgray}边际变化：首次实现上半年盈利，但主要依赖一次性诉讼和解收入；市场关注点从盈利转正转向下半年海外BLA申报进展。\par}
\paragraph{NOM} 药明合联1H26营收预计38.15亿元（同比+41.2\%），BioDlink并表贡献2.25亿元；毛利率稳定在36\%，但运营利润率降至27.4\%（同比-1.1pp），因BioDlink并购管理费用大增77\%。人民币对美元升值4\%导致汇兑损失，净利润仅7.47亿元（同比+0.2\%）。FY26F净利润预测下调22.3\%，目标价从87.73港元调整至82.60港元。
{\small\color{refgray}数据：1H26营收38.15亿元，同比+41.2\%\par}
{\small\color{refgray}数据：BioDlink并表贡献2.25亿元\par}
{\small\color{refgray}数据：毛利率36.0\%（同比持平）\par}
{\small\color{refgray}数据：运营利润率27.4\%（同比-1.1pp）\par}
{\small\color{refgray}边际变化：营收增长逻辑未变，但利润释放节奏被汇率与并购费用打乱；市场关注点从收入增速转向利润端一次性因素消退后的恢复节奏。\par}
\paragraph{NOM} 药明生物1H26营收预计115亿元（同比+15.5\%），早期研究服务（46亿元，+12\%）与商业化项目（51亿元，+18\%）双引擎驱动。毛利率42.0\%（同比-0.7pp），经营利润率27.8\%（-0.9pp），受人民币升值与投资亏损影响，归母净利润约20亿元（同比-15.3\%）。下半年非经常性项目影响减弱，2H26净利润预计29亿元（同比+14\%）。
{\small\color{refgray}数据：1H26营收115亿元，同比+15.5\%\par}
{\small\color{refgray}数据：早期研究服务收入46亿元，同比+12\%\par}
{\small\color{refgray}数据：商业化项目收入51亿元，同比+18\%\par}
{\small\color{refgray}数据：前三大商业化项目2026E各贡献超1亿美元收入\par}
{\small\color{refgray}边际变化：收入端双位数增长，但利润端受汇率与投资亏损双重挤压；市场关注点从短期利润扰动转向下半年非经常性项目影响减弱后的恢复。\par}
\paragraph{GS} 辉瑞2Q26营收预计145.61亿美元（略高于市场共识143.78亿美元），Paxlovid收入1.45亿美元（低于共识10\%），新冠处方量同比下降77\%、环比下降72\%。估值处于历史低位（2027年PE 8.75倍，同业均值14.5倍），但市场焦点已转向下半年两项关键临床数据：mevrometostat 3期（7月）与GLP-1/Amylin早期数据。CFO离职与管线挫折压制情绪，股息率7\%但支付率已超自由现金流。
{\small\color{refgray}数据：2Q26营收预计145.61亿美元（共识143.78亿美元）\par}
{\small\color{refgray}数据：Paxlovid收入1.45亿美元（低于共识10\%）\par}
{\small\color{refgray}数据：新冠处方量同比-77\%，环比-72\%\par}
{\small\color{refgray}数据：2027年PE 8.75倍（同业均值14.5倍）\par}
{\small\color{refgray}边际变化：市场关注点从季度业绩转向下半年关键临床数据；估值处于历史低位，但催化剂不确定性高，市场情绪谨慎。\par}
\paragraph{JPM} 康方生物HARMONi试验更新OS分析：西方患者OS HR收敛至0.76，与全球ITT人群（0.76）及中国HARMONi-A试验结果一致，缓解了OS获益偏亚洲的担忧。但此为计划外分析，无正式p值，仅名义p值显示有利趋势，不能据此宣称统计学显著性。BLA审评（PDUFA 2026年11月14日）仍将依赖预设OS分析（HR 0.79，p=0.057）与PFS数据（HR 0.52，p<0.00001）。
{\small\color{refgray}数据：西方患者OS HR 0.76（与ITT人群一致）\par}
{\small\color{refgray}数据：ITT人群OS HR：0.79→0.78→0.76（2025年4月/9月/2026年6月）\par}
{\small\color{refgray}数据：预设OS分析HR 0.79，p=0.057（未达统计学显著性）\par}
{\small\color{refgray}数据：PFS HR 0.52，p<0.00001\par}
{\small\color{refgray}边际变化：非预设分析显示OS HR持续下降，东西方数据趋同，缓解了市场对OS获益偏亚洲的担忧；但分析性质为探索性，不改变BLA审评的核心证据基础。\par}
\paragraph{NOM} 康方生物1H26F总收入预计24亿元（同比+68\%），药品销售22亿元（同比+55\%，环比+33\%），卡度尼利单抗与依沃西单抗持续放量；合作收入2亿元（伊努西单抗授权）。毛利率提升2.3pp至81.7\%，运营费用增速（11\%）远低于收入增速，营业利润从亏损4.14亿元转为盈利2.3亿元（营业利润率9.7\%），净利润约4600万元（去年同期亏损5.7亿元）。FY26F净利润预测下调79.4\%，反映合作收入减少与运营费用上调。
{\small\color{refgray}数据：1H26F总收入24亿元，同比+68\%\par}
{\small\color{refgray}数据：药品销售22亿元，同比+55\%，环比+33\%\par}
{\small\color{refgray}数据：合作收入2亿元（伊努西单抗授权）\par}
{\small\color{refgray}数据：毛利率81.7\%（同比+2.3pp）\par}
{\small\color{refgray}边际变化：首次实现上半年盈利，核心产品放量驱动经营利润转正；但FY26F净利润预测大幅下调，反映合作收入确认节奏慢于预期。\par}
\paragraph{NOM} 恒瑞医药2Q26总收入预计86.29亿元（同比+1\%），药品销售78亿元（同比+10\%），创新药放量覆盖仿制药下滑；合作收入约8亿元（去年同期14亿元高基数）。归母净利润29亿元（同比+11\%），创单季新高，利润弹性来自销售费用率下降与海外子公司估值变动收益。FY26收入/净利润预测下调7.8\%/18.4\%，因对外授权收入确认慢于预期。
{\small\color{refgray}数据：2Q26总收入86.29亿元，同比+1\%\par}
{\small\color{refgray}数据：药品销售78亿元，同比+10\%\par}
{\small\color{refgray}数据：合作收入约8亿元（去年同期14亿元）\par}
{\small\color{refgray}数据：毛利率85.9\%（同比-1.1pp）\par}
{\small\color{refgray}边际变化：利润创新高但收入增速仅1\%，利润弹性来自费用压缩与一次性收益；市场关注点从收入规模转向利润质量，但对外授权收入确认节奏成为全年业绩关键变量。\par}
\subsubsection*{关键数据}
\begin{itemize}
  \item 科伦博泰1H26净利润7.31亿元（去年同期亏损1.45亿元），诉讼和解收入7.5亿元
  \item 药明合联1H26净利润7.47亿元（同比+0.2\%），人民币对美元升值4\%
  \item 药明生物1H26归母净利润20亿元（同比-15.3\%），前三大商业化项目2026E各贡献超1亿美元收入
  \item 辉瑞2Q26营收预计145.61亿美元，2027年PE 8.75倍（同业均值14.5倍）
  \item 康方生物HARMONi试验西方患者OS HR 0.76，PDUFA日期2026年11月14日
  \item 康方生物1H26F净利润4600万元（去年同期亏损5.7亿元），营业利润率9.7\%
  \item 恒瑞医药2Q26归母净利润29亿元（同比+11\%），单季新高
  \item 恒瑞医药FY26收入/净利润预测下调7.8\%/18.4\%
\end{itemize}
{\small\color{refgray}本节综合 7 篇研究；涉及 NOM、GS、JPM。}
\newpage
\subsection{消费与汽车：结构分化与AI驱动新叙事}
\textbf{消费与汽车板块呈现显著的结构性分化：啤酒行业总量趋平但企业间份额加速转移，华润啤酒表现突出；汽车零部件领域，潍柴动力的估值逻辑正与AI硬件主题挂钩，而恒立液压的滚珠丝杠订单远超指引，机器人业务毛利率有望高于均值。}
\subsubsection*{主流共识}
\begin{itemize}
  \item 啤酒行业整体终端价值增速放缓至+1\%，但企业间分化加剧，华润啤酒增速提升至+6\%，百威中国持续下滑至-4\%。
  \item 汽车芯片行业复苏范围扩大，德州仪器汽车业务实现同比中双位数增长，工业业务保持约30\%同比增长。
  \item 双环传动新能源车齿轮与工程机械齿轮是2Q26F收入增长核心，内燃机齿轮降幅收窄至个位数。
\end{itemize}
\subsubsection*{分歧与条件差异}
\begin{itemize}
  \item 潍柴动力报价与AI硬件指数高度相关（R²=0.84），而中国重汽报价仍与出口增速滞后相关（R²=0.22），两者估值框架已分化。
  \item 恒立液压滚珠丝杠业务年内订单已达5亿元，远超全年3亿元指引，但毛利率受产品组合影响短期承压。
  \item 德州仪器汽车业务回暖，但定价贡献有限，增长主要来自出货量，且高库存导致资本占用较高。
\end{itemize}
\subsubsection*{机构观点}
\paragraph{Bernstein} 6月中国啤酒行业整体终端价值增速从+2\%放缓至+1\%，餐饮渠道增速从+3\%降至+1\%，零售渠道从+1\%升至+2\%。华润啤酒增速从+3\%提升至+6\%，零售渠道反弹至+8\%，广东餐饮增速达+46\%，中高端产品仍是核心增长引擎。百威中国降幅扩大至-4\%，餐饮渠道从+2\%回落至-1\%，复苏仍局限于少数区域。燕京啤酒（未覆盖）以+15\%领跑，青岛啤酒降幅收窄至-1\%。行业总量增长趋平，但份额转移加速。
{\small\color{refgray}数据：6月啤酒行业整体价值增速+1\%（5月+2\%）\par}
{\small\color{refgray}数据：华润啤酒6月价值增速+6\%（5月+3\%）\par}
{\small\color{refgray}数据：百威中国6月价值增速-4\%（5月-3\%）\par}
{\small\color{refgray}数据：燕京啤酒6月价值增速+15\%\par}
{\small\color{refgray}边际变化：华润啤酒零售渠道从5月的+1\%反弹至6月的+8\%，此前该渠道长期低迷；百威中国降幅扩大，复苏缺乏系统性支撑。\par}
\paragraph{Citi} 潍柴动力与中国重汽的报价驱动力已结构性分化。Citi构建的AIDC/SOFC/AI硬件加权指数（Generac与Bloom Energy各50\%权重）与潍柴报价今年以来的R²达0.84，而与重汽的R²仅0.22。潍柴的估值叙事已部分脱离传统重卡周期，与全球AI数据中心电源需求挂钩。重汽报价仍与出口增速存在约15天滞后关系，7月出口走强可能支撑其报价至8月中旬。
{\small\color{refgray}数据：潍柴报价与AI硬件指数R²=0.84\par}
{\small\color{refgray}数据：重汽报价与AI硬件指数R²=0.22\par}
{\small\color{refgray}数据：重汽报价与出口增速存在15天滞后\par}
{\small\color{refgray}数据：Generac和Bloom Energy各占指数50\%权重\par}
{\small\color{refgray}边际变化：潍柴报价从传统重卡周期驱动转向AI硬件主题驱动，市场定价逻辑已发生根本变化。\par}
\paragraph{Citi} 恒立液压滚珠丝杠和直线导轨业务年内累计订单已达5亿元，远超公司全年3亿元收入指引，现有订单已覆盖指引的167\%。日本THK和NSK在中国市场交货周期延长至9个月，为恒立创造进口替代窗口。增量来自注塑机、机床等传统领域国产替代，以及AI和人形机器人产业对机床需求的拉动。宁德时代换电站建设也可能带来额外需求，单个换电站价值量约5万元。新剑传动IPO显示，物理AI线性执行器2023-2025年毛利率分别为62.4\%、70.3\%和59.5\%，恒立机器人行星滚柱丝杠毛利率有望高于公司整体均值。
{\small\color{refgray}数据：滚珠丝杠年内订单5亿元，全年指引3亿元\par}
{\small\color{refgray}数据：现有订单覆盖指引167\%\par}
{\small\color{refgray}数据：日本THK/NSK交货周期延长至9个月\par}
{\small\color{refgray}数据：新剑传动AI执行器毛利率59.5\%-70.3\%\par}
{\small\color{refgray}边际变化：订单远超指引，进口替代窗口因日本同行产能紧张而扩大；人形机器人部件毛利率参照新剑传动数据，可能成为增量贡献。\par}
\paragraph{NOM} 双环传动2Q26F营收延续增长，新能源车齿轮同比增速10\%-20\%，4月持平但5-6月转正，同轴齿轮6月开始放量。内燃机齿轮降幅从1Q26的20\%-30\%收窄至个位数，下半年有望进一步收窄。工程机械齿轮维持20\%-30\%同比增速，特斯拉Semi开始贡献少量收入。机器人减速器单机价值约6000元（6个RV减速器，单价约1000元），产能扩张已提上日程。但2Q26F毛利率不太可能超过1Q26，受产品组合变化影响。
{\small\color{refgray}数据：新能源车齿轮同比增速10\%-20\%\par}
{\small\color{refgray}数据：内燃机齿轮降幅收窄至个位数（1Q26为20\%-30\%）\par}
{\small\color{refgray}数据：工程机械齿轮同比增速20\%-30\%\par}
{\small\color{refgray}数据：机器人减速器单机价值约6000元\par}
{\small\color{refgray}边际变化：收入结构从传统内燃机向新能源和工程机械切换，机器人减速器进入客户交付和产能扩张早期阶段。\par}
\paragraph{MS} 德州仪器二季度营收54.63亿美元，环比增长13.2\%，同比增长22.8\%，毛利率回升至61.4\%。汽车业务出现实质性回暖，同比中双位数增长、环比高个位数增长，管理层称“需求无处不在”。工业业务同比增长约30\%、环比约10\%，仍比2022年峰值低5-6个百分点。三季度指引中点59亿美元，环比增长约8\%，增长主要来自出货量，定价贡献几乎可忽略。2026年资本开支指引偏向上限（20-30亿美元），MS将预测从26亿上调至28亿，自由现金流改善时间点可能晚于预期。
{\small\color{refgray}数据：2Q26营收54.63亿美元，环比+13.2\%，同比+22.8\%\par}
{\small\color{refgray}数据：毛利率61.4\%（高于预期1.4个百分点）\par}
{\small\color{refgray}数据：汽车业务同比中双位数增长\par}
{\small\color{refgray}数据：工业业务同比约30\%，环比约10\%\par}
{\small\color{refgray}边际变化：汽车业务从过去几个季度的低迷中实质性回暖，成为复苏扩大的关键信号；但增长主要靠出货量而非定价，且高库存导致资本占用较高。\par}
\subsubsection*{关键数据}
\begin{itemize}
  \item 6月啤酒行业整体价值增速+1\%（5月+2\%）
  \item 华润啤酒6月价值增速+6\%（5月+3\%）
  \item 潍柴报价与AI硬件指数R²=0.84
  \item 恒立液压滚珠丝杠年内订单5亿元，全年指引3亿元
  \item 新剑传动AI执行器毛利率59.5\%-70.3\%
  \item 双环传动新能源车齿轮同比增速10\%-20\%
  \item 德州仪器2Q26营收54.63亿美元，汽车业务同比中双位数增长
  \item 德州仪器3Q26指引中点59亿美元，环比+8\%
\end{itemize}
\begin{figure}[htbp]
\centering
\includegraphics[width=0.88\linewidth]{figures/fig_041.jpg}
\caption{Figure 1 - Kyle Wu Figure 1. YTD Share Price Movements of Weichai/Sinotruk/Generac/BE © 2026 Citi Inc. No redistribution without Citi's written permission. Note: We develop a AIDC/SOFC/AI-Hardware weighted index, by assigning 50\%/50\% weight on daily share price of Genera}
\end{figure}
\begin{figure}[htbp]
\centering
\includegraphics[width=0.88\linewidth]{figures/fig_058.jpg}
\caption{Figure 1 - Figure 1. Hengli: Ball screw and linear guide 2026 revenue guidance vs. order YTD © 2026 Citi Inc. No redistribution without Citi's written permission. Figure 2. Seenpin: Physical AI linear actuator revenue and GPM, 2023-25 © 2026 Citi Inc. No redistribution w}
\end{figure}
\begin{figure}[htbp]
\centering
\includegraphics[width=0.88\linewidth]{figures/fig_006.jpg}
\caption{EXHIBIT 1 - EXHIBIT 2: We see a 73-96\% correlation between BigOne Lab quarterly brewer value growth and reported revenue growth Regression R-square of BigOne National Sell-out Value YoY \% vs Reported Beer Revenue YoY \% (22Q3-26Q1) CRBeer rep}
\end{figure}
\begin{figure}[htbp]
\centering
\includegraphics[width=0.88\linewidth]{figures/fig_139.jpg}
\caption{Exhibit 1 - \#\# OWNERSHIP POSITIONING Refinitiv; MSPB Content. Includes certain hedge fund exposures held with MSPB. Information may be inconsistent with or may not reflect broader market trends. Long/Short Ratio = Long Exposure / Short exposure. Sector \% of Total Net Expo}
\end{figure}
{\small\color{refgray}本节综合 5 篇研究；涉及 Bernstein、Citi、NOM、MS。}
\newpage
\subsection{工业与机器人：结构性分化与成本节奏错配}
\textbf{全球工业与机器人领域呈现结构性分化：机器人手术渗透率加速提升但竞争格局松动，电梯行业订单稳健但成本通胀侵蚀利润，中国钢铁供需背离、库存被动累积。}
\subsubsection*{主流共识}
\begin{itemize}
  \item 全球软组织机器人手术渗透率未来五年将从33\%升至48\%，年复合增速约11\%，增长来自开放手术转换和腹腔镜升级。
  \item 电梯行业新安装订单区域分化，亚太（除中国）驱动增长，现代化改造订单保持韧性。
  \item 中国钢铁市场表观消费走弱，产量仍惯性扩张，贸易商库存同比高企。
\end{itemize}
\subsubsection*{分歧与条件差异}
\begin{itemize}
  \item 机器人手术：美国市场渗透率增幅温和（64\%→69\%），增量几乎全部来自开放手术转换；而全球市场因基数低、医生采纳加速，实际增速可能高于调查值。
  \item 电梯行业：Schindler利润率扩张节奏放缓（Q2 adj EBIT +35bps vs Q1），管理层归因于成本通胀加速；但订单积压利润率环比改善，下半年成本压力可能集中释放。
  \item 中国钢铁：长材消费同比降9.1\%，扁平材降2.9\%，但产量仍环比增长；贸易商库存同比+27.1\%，钢厂库存仅+7.6\%，渠道被动累库信号明确。
\end{itemize}
\subsubsection*{机构观点}
\paragraph{DB} 全球机器人手术渗透率加速提升，但竞争格局正在松动。调查显示，某系统全球装机份额预计从81\%降至68\%，美敦力Hugo将成第二大品牌（10\%份额）。美国市场因成熟，渗透率仅从64\%升至69\%，增量来自开放手术转换。手术向门诊手术中心迁移趋势明确，五年后约三分之一手术将在ASC完成。调查对象为SRS会员（平均17年经验），实际全球增速可能更快。
{\small\color{refgray}数据：全球机器人手术渗透率33\%→48\%（5年）\par}
{\small\color{refgray}数据：年复合增长率约11\%\par}
{\small\color{refgray}数据：某系统全球装机份额81\%→68\%\par}
{\small\color{refgray}数据：美国渗透率64\%→69\%\par}
{\small\color{refgray}边际变化：竞争格局从一家独大转向多品牌竞争，且手术地点从医院向ASC迁移加速。\par}
\paragraph{DB} Schindler二季度本地货币订单增长约3\%，与一季度持平，但利润率扩张放缓至+35bps。管理层将原因归结为成本通胀加速，全年成本通胀指引上调至约3500万瑞士法郎，其中约三分之二预计发生在下半年。新安装订单量低两位数增长，由亚太（除中国）驱动；中国区新安装继续调整，变化幅度快于同行。现代化改造订单在高基数上仍增11\%，订单积压利润率环比改善。
{\small\color{refgray}数据：Q2本地货币订单增速约3\%\par}
{\small\color{refgray}数据：adj EBIT利润率同比+35bps（Q1更高）\par}
{\small\color{refgray}数据：全年成本通胀约3500万瑞郎，下半年占2/3\par}
{\small\color{refgray}数据：现代化改造订单同比+11\%（基数+24\%）\par}
{\small\color{refgray}边际变化：利润率扩张节奏放缓，成本通胀压力从上半年延迟至下半年集中体现。\par}
\paragraph{MS} 中国钢铁供需背离加剧。长材表观消费环比-3.8\%、同比-9.1\%，扁平材环比-1.6\%、同比-2.9\%；但长材产量环比+2.5\%，扁平材+0.9\%。贸易商库存环比+0.8\%、同比+27.1\%，钢厂库存持平。高炉利用率89.8\%（环比-0.8pp），电炉开工率55.7\%（环比-1.3pp）。铁矿石港口库存环比-1.8\%，澳巴发运合计减少59万吨。
{\small\color{refgray}数据：长材表观消费同比-9.1\%\par}
{\small\color{refgray}数据：长材产量环比+2.5\%\par}
{\small\color{refgray}数据：贸易商库存同比+27.1\%\par}
{\small\color{refgray}数据：高炉利用率89.8\%（环比-0.8pp）\par}
{\small\color{refgray}边际变化：需求走弱而供给惯性扩张，渠道库存被动累积；电炉开工率快速下降反映短流程利润承压。\par}
\subsubsection*{关键数据}
\begin{itemize}
  \item 全球机器人手术渗透率33\%→48\%（5年）
  \item 某系统全球装机份额81\%→68\%
  \item Schindler Q2 adj EBIT利润率+35bps
  \item 全年成本通胀约3500万瑞郎，下半年占2/3
  \item 长材表观消费同比-9.1\%
  \item 贸易商库存同比+27.1\%
  \item 电炉开工率55.7\%（环比-1.3pp）
  \item 澳巴铁矿石发运减少59万吨
\end{itemize}
\begin{figure}[htbp]
\centering
\includegraphics[width=0.88\linewidth]{figures/fig_135.jpg}
\caption{Exhibit 1 - Exhibit 2: Weekly steel demand Rachel L Zhang Hannah Yang, CFA}
\end{figure}
\begin{figure}[htbp]
\centering
\includegraphics[width=0.88\linewidth]{figures/fig_064.jpg}
\caption{Figure 3 - Figure 3: Adj. EBIT and margin}
\end{figure}
\begin{figure}[htbp]
\centering
\includegraphics[width=0.88\linewidth]{figures/fig_061.jpg}
\caption{Figure 1 - Figure 1: Sales and organic sales growth}
\end{figure}
\begin{figure}[htbp]
\centering
\includegraphics[width=0.88\linewidth]{figures/fig_062.jpg}
\caption{Figure 1 - Figure 1: Sales and organic sales growth}
\end{figure}
{\small\color{refgray}本节综合 3 篇研究；涉及 DB、MS。}
\newpage
\subsection{能源与电动化：欧洲需求回升与供给瓶颈并存}
\textbf{欧洲电动化进程在2026年上半年显著加速，BEV注册量同比增35\%，但充电基础设施季度建设速度同比下降20\%，供需缺口扩大。中国电池龙头宁德时代以创纪录业绩和回购计划强化市场信心，而敏实集团等供应商受益于欧洲需求复苏。}
\subsubsection*{主流共识}
\begin{itemize}
  \item 欧洲BEV市场在2026年上半年实现强劲增长，1-6月注册量同比增35.1\%，渗透率达22.2\%。
  \item 充电需求正从存量利用率提升转向增量车辆入场双重驱动，Q2 BEV注册量增速超整体汽车市场50个百分点。
  \item 宁德时代二季度净利润225.5亿元同比增36\%，销量维持约60\%同比增长，储能占比升至25\%。
\end{itemize}
\subsubsection*{分歧与条件差异}
\begin{itemize}
  \item 充电桩建设速度放缓（季度新增同比降20\%）与需求增长形成矛盾，各国分化明显，德国增速落后于欧盟整体。
  \item 宁德时代单位净利润环比下降7\%至97元/kWh，但EBIT利润率超预期，费用控制与外汇对冲效果存在分歧。
  \item 敏实集团与欧洲电动化趋势高度相关，但铝价上涨对毛利率的影响存在不确定性，历史数据显示韧性。
\end{itemize}
\subsubsection*{机构观点}
\paragraph{DB} 欧洲电动化趋势与敏实集团高度相关：2026年上半年欧洲BEV注册量同比增35.1\%，敏实电池盒业务收入预计同比增33\%至48亿元，推动集团总收入增9\%至134亿元。公司历史毛利率在铝价上涨周期中从32.6\%提升至34\%，显示成本传导能力。上半年净利润预计14.3亿元同比增12\%，市场预期从41.8上调至42.3。
{\small\color{refgray}数据：欧洲1-6月BEV注册量160.8万辆，同比+35.1\%\par}
{\small\color{refgray}数据：敏实电池盒收入上半年预计48亿元，同比+33\%\par}
{\small\color{refgray}数据：敏实上半年净利润预计14.3亿元，同比+12\%\par}
{\small\color{refgray}数据：敏实报价从2024年8月10港元升至2026年2月46.76港元\par}
{\small\color{refgray}边际变化：欧洲电动化从2023-2024年减速转向2024Q4以来恢复增长，敏实报价同步反弹，市场预期上调。\par}
\paragraph{DB} 欧洲充电需求出现非季节性走强，Q2 BEV注册量增速超整体汽车市场50个百分点，但公共充电桩季度建设量仅4.5万个，同比下降20\%，供需缺口扩大。截至Q2末欧盟充电桩总量117万个，同比增18\%，但按当前速度2026年底仅达126.7万个。各国分化明显，意大利同比增21\%领先，德国仅增14\%落后。
{\small\color{refgray}数据：Q2欧盟公共充电桩季度新增4.5万个，同比-20\%\par}
{\small\color{refgray}数据：欧盟充电桩总量117万个，同比+18\%\par}
{\small\color{refgray}数据：德国充电桩增速14\%，低于欧盟整体\par}
{\small\color{refgray}数据：意大利充电桩增速21\%，领先五大市场\par}
{\small\color{refgray}边际变化：充电需求从存量利用率提升转向增量车辆入场双重驱动，但建设速度放缓，供需缺口扩大。\par}
\paragraph{GS} 宁德时代二季度净利润225.5亿元同比增36\%，符合预期；电池销量约235GWh同比增60\%，储能占比从20\%升至25\%。单位净利润97元/kWh环比降7\%，但EBIT利润率17.2\%超预期，费用控制与外汇对冲有效。公司宣布200-400亿元回购注销计划，为A股最大规模，管理层对下半年及2027年需求乐观。
{\small\color{refgray}数据：宁德时代二季度净利润225.5亿元，同比+36\%\par}
{\small\color{refgray}数据：二季度电池销量约235GWh，同比+60\%\par}
{\small\color{refgray}数据：储能占比从20\%升至25\%\par}
{\small\color{refgray}数据：单位净利润97元/kWh，环比-7\%\par}
{\small\color{refgray}边际变化：宁德时代从单纯产能扩张转向兼顾股东回报，大规模回购注销提升每股收益，储能与海外成为新增长引擎。\par}
\subsubsection*{关键数据}
\begin{itemize}
  \item 欧洲1-6月BEV注册量160.8万辆，同比+35.1\%
  \item Q2欧盟公共充电桩季度新增4.5万个，同比-20\%
  \item 宁德时代二季度净利润225.5亿元，同比+36\%
  \item 宁德时代二季度电池销量约235GWh，同比+60\%
  \item 敏实电池盒收入上半年预计48亿元，同比+33\%
  \item 宁德时代单位净利润97元/kWh，环比-7\%
  \item 欧盟充电桩总量117万个，同比+18\%
  \item 宁德时代回购计划200-400亿元，A股最大规模
\end{itemize}
\begin{figure}[htbp]
\centering
\includegraphics[width=0.88\linewidth]{figures/fig_066.jpg}
\caption{Figure 1 - Figure 1: In Europe, the public charging network growth is ongoing at a lower pace y/y...}
\end{figure}
\begin{figure}[htbp]
\centering
\includegraphics[width=0.88\linewidth]{figures/fig_076.jpg}
\caption{Exhibit 1 - Exhibit 1: CATL saw c.60\% yoy volume growth in 2Q26, on doubled yoy ESS and solid power battery CATL quarterly sales volume and yoy chg\% Exhibit 2: 2Q26 recorded unit net profit of Rmb97/kWh, down -7\% qoq from Rmb104/kWh in 1Q26}
\end{figure}
\begin{figure}[htbp]
\centering
\includegraphics[width=0.88\linewidth]{figures/fig_067.jpg}
\caption{Figure 1 - Figure 4: ...necessitating a 7x increase in the build rate despite a lower target by 2030}
\end{figure}
\begin{figure}[htbp]
\centering
\includegraphics[width=0.88\linewidth]{figures/fig_068.jpg}
\caption{Figure 2 - Figure 5: The top-five countries are showing different growth paths, with Italy growing the most strongly...}
\end{figure}
{\small\color{refgray}本节综合 3 篇研究；涉及 DB、GS。}
\newpage
\subsection{金融与银行：香港银行净息差预期修正与股东回报关注}
\textbf{市场对香港银行净息差走势的预期出现关键修正，从“逐年收窄”转向“持稳”，叠加信用成本下降和股东回报增强预期，中银香港等标的的盈利预测和估值逻辑面临重估。}
\subsubsection*{主流共识}
\begin{itemize}
  \item 香港银行净息差此前被普遍认为已见顶，未来将逐年收窄。
  \item 香港商业地产贷款风险是市场关注焦点，但2026年上半年未出现新的实质性违约。
  \item 中银香港具备提升股东回报的能力，市场已部分定价回购或特别股息。
\end{itemize}
\subsubsection*{分歧与条件差异}
\begin{itemize}
  \item 净息差走势：JPM认为中银香港净息差（含掉期收入）在FY26-FY27可维持在1.6\%左右，优于此前每年下降3-4个基点的预测；而市场主流观点仍认为息差将逐步收窄。
  \item 盈利增长驱动：JPM上调FY26和FY27盈利预测8\%和9\%，主要基于净息差持稳和信用成本下降；但部分机构认为净利息收入增长仍面临利率下行压力。
  \item 股东回报力度：JPM基准情景预计每年特别股息0.95港元，额外贡献约2个百分点总回报率；但市场对实际回购或派息规模存在分歧。
\end{itemize}
\subsubsection*{机构观点}
\paragraph{JPM} JPM调整中银香港净息差预期，从“逐年收窄3-4个基点”修正为FY26-FY27均维持在约1.6\%（含掉期收入调整），主要基于美联储偏鹰表态及HIBOR远期曲线上升。这一修正带动FY26和FY27盈利预测分别上调8\%和9\%，净利息收入在FY26有望实现高个位数增长。同时，香港商业地产贷款组合在2026年上半年未出现新的实质性违约，信用成本同比若延续下降趋势将构成额外盈利支撑。股东回报方面，中银香港在FY26-FY28期间有能力每年提供最高约3个百分点的额外回报，基准情景为每年特别股息0.95港元，额外贡献约2个百分点总回报率。估值对应FY27预期市净率1.4倍，基于约13\%的净资产回报率预测。
{\small\color{refgray}数据：净息差预期从逐年收窄3-4个基点修正为FY26-FY27维持约1.6\%\par}
{\small\color{refgray}数据：FY26和FY27盈利预测分别上调8\%和9\%\par}
{\small\color{refgray}数据：FY26净利息收入有望实现高个位数增长\par}
{\small\color{refgray}数据：2026年上半年香港商业地产贷款未出现新的实质性违约\par}
{\small\color{refgray}边际变化：净息差预期从“逐年收窄”转为“持稳”，是本次报告最关键的判断变化，直接带动盈利预测上修和估值重估。\par}
\subsubsection*{关键数据}
\begin{itemize}
  \item 净息差预期从逐年收窄3-4个基点修正为FY26-FY27维持约1.6\%
  \item FY26和FY27盈利预测分别上调8\%和9\%
  \item FY26净利息收入有望实现高个位数增长
  \item 2026年上半年香港商业地产贷款未出现新的实质性违约
  \item 基准情景每年特别股息0.95港元，额外贡献约2个百分点总回报率
  \item FY26-FY28每年最高可提供约3个百分点额外股东回报
  \item 对应FY27预期市净率1.4倍，净资产回报率约13\%
  \item 年初至今股价跑赢恒生指数约29个百分点
\end{itemize}
{\small\color{refgray}本节综合 1 篇研究；涉及 JPM。}
\newpage
\subsection{区域市场与主题：AI驱动增长、能源冲击与政策应对}
\textbf{全球增长韧性超预期，AI资本开支成为结构性支撑，但能源价格上行与利率路径不确定性正在重塑市场定价；中国市场呈现生产端修复与内需偏弱的分化，政策托底信号明确但效果待观察。}
\subsubsection*{主流共识}
\begin{itemize}
  \item AI资本开支是全球制造业PMI回升的核心驱动力，美国AI投资已从2025年不足3000亿美元跃升至2026年Q1的4300亿美元，全年或达6000亿美元（接近GDP的2\%），并向亚洲扩散。
  \item 全球增长对石油依赖度降低，Citi维持2026年全球增长预测2.5\%附近，仅较年初下调零点几个百分点；能源冲击更多通过通胀预期传导而非直接压制增长。
  \item 中国高端美妆市场2025年增长6\%至1860亿元，雅诗兰黛份额从17.6\%回升至18.7\%，SK-II从2.9\%反弹至3.3\%，但增长节奏不规律，可持续性存疑。
\end{itemize}
\subsubsection*{分歧与条件差异}
\begin{itemize}
  \item 能源价格对利率路径的影响：GS认为若美联储7月按兵不动且沟通有限，长端通胀不确定性溢价将重估；而Citi认为全球通胀预期已上调约0.75个百分点，主要央行政策路径已转向审慎。
  \item 中国内需修复节奏：GS的CAI显示6月环比年化+5.3\%（5月+4.2\%），但进口隐含需求指标偏弱，生产端修复与需求端调整的分化持续；JPM认为二季度盈利预期可实现，但板块间分化显著。
  \item 美国中期选举对市场的影响：GS认为选举结果本身不是波动主因，但隐含相关性处于历史低位叠加宏观不确定性上升，指数波动率可能被低估；历史数据显示8月至选举日标普500中位回报为零。
\end{itemize}
\subsubsection*{机构观点}
\paragraph{Bernstein} 地平线可转债重组消除4.9\%潜在稀释，实际效果相当于1\%股份回购；大众锁定期延长12个月，零息新可转债节省约2000万美元利息。大众L3/L4合作验证地平线全栈技术价值，首批L3交付目标2027年下半年。市场当天3.8\%波动属误判，两项变化均正面。
{\small\color{refgray}数据：潜在稀释从4.9\%降至3.9\%，降低1个百分点\par}
{\small\color{refgray}数据：新可转债转股价5.55港元，较收盘价溢价16.8\%\par}
{\small\color{refgray}数据：节省约2000万美元名义利息支出\par}
{\small\color{refgray}数据：大众持股锁定期延长12个月\par}
{\small\color{refgray}边际变化：从市场误判为负面到确认资本结构优化与战略合作升级，稀释风险降低且融资成本改善。\par}
\paragraph{Bernstein} 雅诗兰黛在中国高端美妆市场的反弹并非来自国产品牌让出份额，而是自身品牌投入驱动。2025年中国高端美妆市场增长6\%至1860亿元，雅诗兰黛份额从17.6\%回升至18.7\%，SK-II从2.9\%反弹至3.3\%。但增长节奏不规律，2019-2025年CAGR仅5.1\%，远低于疫情前双位数增速，可持续性存疑。
{\small\color{refgray}数据：2025年中国高端美妆市场规模1860亿元（约270亿美元）\par}
{\small\color{refgray}数据：2019-2025年CAGR 5.1\%\par}
{\small\color{refgray}数据：雅诗兰黛份额从17.6\%升至18.7\%\par}
{\small\color{refgray}数据：SK-II份额从2.9\%升至3.3\%\par}
{\small\color{refgray}边际变化：从市场核心争论'能否复苏'转向'复苏能否持续'，数据不支持国产品牌退让的流行叙事。\par}
\paragraph{Bernstein} NextEra Energy Q2调整后EPS 1.15美元，高于预期的1.13美元，同比增长9.5\%。FPL大型负荷需求管线从6GW更新至8GW（2032年预期），另有12GW深入讨论中。NEER季度新增3.6GW可再生能源+储能订单，总积压订单从33GW增至35.1GW，超500MW项目重新签约溢价约20美元/MWh。数据中心需求验证逻辑，30个潜在枢纽洽谈中。
{\small\color{refgray}数据：Q2调整后EPS 1.15美元，同比增长9.5\%\par}
{\small\color{refgray}数据：FPL大型负荷需求8GW纳入2032年预期（此前6GW）\par}
{\small\color{refgray}数据：NEER季度新增3.6GW订单，总积压35.1GW\par}
{\small\color{refgray}数据：重新签约溢价约20美元/MWh\par}
{\small\color{refgray}边际变化：FPL需求管线从6GW更新至8GW，基于实际谈判进展而非预测，数据中心需求兑现速度加快。\par}
\paragraph{Bernstein} 全球供应链运价持续上涨，海运运价较冲突前上涨超100\%，空运运价同比上升25-30\%。5月海运集装箱量同比增长4\%，马士基上调全年行业增长预期至约4\%。美国现货整车运价（剔除燃油）同比上涨48\%，远高于5月的34\%，供给端收紧是主因（监管执行加强+燃油成本上升）。不同运输模式和区域市场出现分化。
{\small\color{refgray}数据：海运运价较冲突前上涨超100\%\par}
{\small\color{refgray}数据：空运运价同比上升25-30\%\par}
{\small\color{refgray}数据：美国现货整车运价同比上涨48\%（5月为34\%）\par}
{\small\color{refgray}数据：5月海运集装箱量同比增长4\%\par}
{\small\color{refgray}边际变化：运价上涨从单一地缘因素驱动转向需求脉冲与运力结构性收紧共同支撑，美国公路运价涨幅加速。\par}
\paragraph{NOM} VNET Q2收入预计27.6亿元，同比增长13.2\%，批发IDC收入同比增长36.4\%为核心引擎，但零售IDC仅增1.2\%。新增电力爬坡受国内芯片产能限制而非需求不足。调整后EBITDA利润率同比扩张2.6个百分点至32.7\%，但环比微降0.4个百分点。2026-2028年收入预测下调1.2-2.7\%，EBITDA预测上调0.6-1.9\%。
{\small\color{refgray}数据：Q2收入预计27.6亿元，同比增长13.2\%\par}
{\small\color{refgray}数据：批发IDC收入同比增长36.4\%\par}
{\small\color{refgray}数据：调整后EBITDA利润率32.7\%，同比+2.6pp\par}
{\small\color{refgray}数据：2026-2028年收入预测下调1.2-2.7\%\par}
{\small\color{refgray}边际变化：收入增速受制于芯片产能爬坡节奏，但批发业务放量与成本控制推动EBITDA利润率改善，定价能力预计2027-2028年稳中有升。\par}
\paragraph{Citi} 全球增长韧性超预期，2026年增长预测维持2.5\%附近，仅较年初下调零点几个百分点。AI投资从2025年不足3000亿美元跃升至2026年Q1的4300亿美元，全年或达6000亿美元（接近GDP的2\%），并向亚洲扩散，成为全球制造业PMI回升的核心驱动力。全球通胀预期较2月上调约0.75个百分点，主要央行政策路径转向审慎。
{\small\color{refgray}数据：2026年全球增长预测2.5\%附近\par}
{\small\color{refgray}数据：美国AI投资从<3000亿美元（2025）升至Q1 4300亿美元\par}
{\small\color{refgray}数据：全年AI投资或达6000亿美元（GDP的2\%）\par}
{\small\color{refgray}数据：全球通胀预期较2月上调约0.75个百分点\par}
{\small\color{refgray}边际变化：从市场担忧能源冲击导致显著下滑，到确认AI投资成为结构性增量，增长韧性超预期。\par}
\paragraph{GS} 雪佛龙和康菲石油因多元化高利润率资产组合和成本削减路径被视为更具吸引力。雪佛龙预计今年产量增长7-10\%，在70美元/桶布伦特假设下自由现金流和EPS年均增长约10\%。康菲石油到2029年自由现金流增长约70亿美元（按70美元WTI），每股自由现金流CAGR 20-25\%。埃克森美孚EV/DACF约9.5倍，估值相对较高。
{\small\color{refgray}数据：雪佛龙产量增长7-10\%，自由现金流年均增长约10\%\par}
{\small\color{refgray}数据：康菲石油到2029年自由现金流增长约70亿美元\par}
{\small\color{refgray}数据：康菲石油每股自由现金流CAGR 20-25\%\par}
{\small\color{refgray}数据：雪佛龙与微软签署20年期天然气发电协议（2.67GW）\par}
{\small\color{refgray}边际变化：石油巨头从单纯油气生产商向'能源+基础设施'角色延伸，雪佛龙与微软的电力协议改变估值框架认知。\par}
\paragraph{GS} 油价自7月CPI公布以来上涨约15\%，推动7月美联储会议加息定价从边缘可能性变为显著不确定性。GS维持7月按兵不动基准判断，但若沟通有限，长端利率和通胀不确定性溢价将面临重估。欧洲前端利率仍与能源价格高度绑定，市场已定价9月加息。英国Gilt风险溢价居高不下，日本曲线变平取决于BoJ表态。
{\small\color{refgray}数据：油价自7月CPI公布以来上涨约15\%\par}
{\small\color{refgray}数据：7月美联储会议加息定价从边缘可能性变为显著不确定性\par}
{\small\color{refgray}数据：欧洲市场已定价9月加息\par}
{\small\color{refgray}数据：英国市场已计入2027年近3次加息\par}
{\small\color{refgray}边际变化：能源价格上行正在重塑全球主要地区利率路径预期，央行沟通方式成为决定市场反应的关键变量。\par}
\paragraph{GS} 中国6月CAI环比年化+5.3\%，较5月+4.2\%进一步回升，由制造业和消费拉动。但进口隐含国内需求指标显示二季度内需增长偏弱，呈现生产端修复与需求端调整的分化。金融条件指数6月收紧，信用脉冲转负。制造业和建筑活动代理指标均环比改善，但投资追踪器显示二季度实际增长偏弱。
{\small\color{refgray}数据：6月CAI环比年化+5.3\%（5月+4.2\%）\par}
{\small\color{refgray}数据：制造业增长代理和建筑增长代理均环比改善\par}
{\small\color{refgray}数据：金融条件指数6月收紧\par}
{\small\color{refgray}数据：信用脉冲已转为负值\par}
{\small\color{refgray}边际变化：从5月+4.2\%回升至6月+5.3\%，但内需仍偏弱，结构上生产端强于需求端。\par}
\paragraph{GS} SAP Q2 CCB同比增长26\%（有机约25\%），高于GS和一致预期的25\%和24\%，驱动因素为Cloud ERP Suite持续强势及Autonomous Suite和Business AI Platform新动能。GS将FY26云收入增速预测从24.0\%上调至24.4\%，但EBIT增长预测从14.8\%下调至14.5\%，反映云毛利率下降、运营支出略高及收购稀释。管理层重申全年指引。
{\small\color{refgray}数据：Q2 CCB同比增长26\%（有机约25\%）\par}
{\small\color{refgray}数据：FY26云收入增速预测从24.0\%上调至24.4\%\par}
{\small\color{refgray}数据：FY26 EBIT增长预测从14.8\%下调至14.5\%\par}
{\small\color{refgray}数据：收购Dremio和Prior Labs带来约1亿欧元稀释\par}
{\small\color{refgray}边际变化：CCB增长验证产品周期韧性，但利润端因研发投入和收购稀释小幅下修，AI变现验证点在未来几个季度。\par}
\paragraph{GS} 美国中期选举前市场进入观望期，1974年以来13次中期选举中，8月初至选举日标普500中位回报为零，选举后三个月中位回报6\%。隐含相关性处于历史低位，指数波动率被低估。共同基金在选举前三个月平均提高现金仓位0.4\%，外国投资者减持0.1\%美国权益资产。10年期实际利率升至2023年以来最高。
{\small\color{refgray}数据：8月初至选举日标普500中位回报0\%\par}
{\small\color{refgray}数据：选举后三个月标普500中位回报6\%\par}
{\small\color{refgray}数据：隐含相关性处于近几十年最低水平\par}
{\small\color{refgray}数据：共同基金选举前三个月现金仓位提高0.4\%\par}
{\small\color{refgray}边际变化：从关注选举结果本身转向关注宏观不确定性上升、低隐含相关性与利率波动叠加对指数波动率的推升作用。\par}
\paragraph{JPM} 中国诚通和国新控股宣布合计600亿元场内公司与ETF购买计划，科技公司被明确纳入。ETF流量追踪显示基准ETF及科创50、科创芯片等AI/科技主题产品资金快速逆转，A股每日回购规模升至年内高点。MSCI中国和沪深300 Q2 EPS共识增速分别为2.8\%和15.7\%，可实现但板块间分化显著：IT和工业扩张，消费和地产调整。互联网板块呈现'收入向上、盈利持平'动态。
{\small\color{refgray}数据：中国诚通和国新控股合计600亿元购买计划\par}
{\small\color{refgray}数据：沪深300 Q2 EPS共识增速15.7\%\par}
{\small\color{refgray}数据：MSCI中国 Q2 EPS共识增速2.8\%\par}
{\small\color{refgray}数据：A股每日回购规模升至年内高点\par}
{\small\color{refgray}边际变化：政策托底信号明确，资金快速逆转，但二季度业绩期成为检验板块真实盈利动能的关键窗口，非AI和国内需求板块基本面支撑有限。\par}
\subsubsection*{关键数据}
\begin{itemize}
  \item 美国AI投资从2025年不足3000亿美元跃升至2026年Q1的4300亿美元，全年或达6000亿美元（GDP的2\%）
  \item 中国6月CAI环比年化+5.3\%（5月+4.2\%），但内需仍偏弱
  \item 海运运价较冲突前上涨超100\%，美国现货整车运价同比上涨48\%
  \item 中国诚通和国新控股合计600亿元场内购买计划，A股回购升至年内高点
  \item SAP Q2 CCB同比增长26\%，FY26云收入增速预测上调至24.4\%
  \item NextEra Energy FPL大型负荷需求从6GW更新至8GW
  \item 雪佛龙与微软签署20年期2.67GW天然气发电协议
  \item 2025年中国高端美妆市场增长6\%至1860亿元，雅诗兰黛份额回升至18.7\%
\end{itemize}
\begin{figure}[htbp]
\centering
\includegraphics[width=0.88\linewidth]{figures/fig_048.jpg}
\caption{Figure 1 - \#\# Renewed Uncertainties—But Continued Resilience Tensions between the United States and Iran continue to cast a long shadow over the global oil market. Only few weeks ago, oil prices were approaching \textbackslash{}\$70/barrel and shipments through the Strait of Hormuz look}
\end{figure}
\begin{figure}[htbp]
\centering
\includegraphics[width=0.88\linewidth]{figures/fig_090.jpg}
\caption{Exhibit 1 - Exhibit 1: Our China Current Activity Indicator (CAI) rose to +5.3\% mom annualized sa in June, vs. +4.2\% in May Exhibit 2: June's improvement in CAI was driven by manufacturing and consumption Exhibit 3: Our proprietary import-}
\end{figure}
\begin{figure}[htbp]
\centering
\includegraphics[width=0.88\linewidth]{figures/fig_105.jpg}
\caption{Exhibit 1 - Exhibit 1: Policy uncertainty typically rises ahead of midterm elections Exhibit 2: S\&P 500 volatility usually rises ahead of midterm elections Alongside elevated uncertainty, mutual funds and foreign investors have typically d}
\end{figure}
\begin{figure}[htbp]
\centering
\includegraphics[width=0.88\linewidth]{figures/fig_115.jpg}
\caption{Figure 1 - Figure 1: Accumulative benchmark ETF inflows since 2024: market tends to recover following the ETF inflow upticks Figure 2: Accumulative STAR50 + STAR Chip ETF inflows since 2024 Figure 3: Accumulative A-share ETF net flow has}
\end{figure}
{\small\color{refgray}本节综合 12 篇研究；涉及 Bernstein、NOM、Citi、GS、JPM。}
\newpage
\subsection{区域市场与主题：情绪分化与个股结构性机会}
\textbf{中国A股短期情绪偏冷但修复窗口仍在，港股因南向资金与科技盈利加速更具吸引力；个股层面，PayPay与第一拖拉机股份均呈现利润增速快于收入的盈利结构改善，而BeOne面临一次性税款导致的利润错位。}
\subsubsection*{主流共识}
\begin{itemize}
  \item 中国A股市场短期情绪指标走弱，政策空窗期与新股发行导致资金观望。
  \item 科技自主、先进制造和能源安全仍是政策优先方向，地产数据持续疲弱。
  \item 个股层面，核心业务增长轨迹未变，但市场对非核心风险（如寿险投资、税款）过度反应。
\end{itemize}
\subsubsection*{分歧与条件差异}
\begin{itemize}
  \item A股vs港股：A股情绪偏冷且成交萎缩，港股则受益于南向资金持续流入和科技股盈利加速，相对吸引力更强。
  \item 政策预期：部分观点认为政策将加速使用剩余财政额度，但不会追加预算；另一些观点期待更广泛的消费刺激。
  \item PayPay估值：市场因寿险投资不确定性给予折价，但MS认为核心业务增长未变，估值已反映悲观预期。
\end{itemize}
\subsubsection*{机构观点}
\paragraph{MS} 中国A股短期情绪偏冷，MSASI加权读数从57\%降至37\%，创业板成交额环比降6\%至6670亿元，融资余额降6\%至2.73万亿元，但30日RSI稳定，非恐慌性抛售。政策侧重“速度而非规模”，约2万亿元预算内财政额度加速落地，但无新增预算。7月底至8月仍处可持续修复窗口，技术硬件与AI板块盈利加速是底层支撑。港股更值得关注：7月至今南向资金净流入约111亿美元，科技股盈利增速加快，电商价格战接近尾声，AI商业化推进。
{\small\color{refgray}数据：MSASI加权读数从57\%降至37\%（7月15日至22日）\par}
{\small\color{refgray}数据：创业板成交额环比降6\%至6670亿元\par}
{\small\color{refgray}数据：融资余额降6\%至2.73万亿元\par}
{\small\color{refgray}数据：南向资金7月至今净流入约111亿美元\par}
{\small\color{refgray}边际变化：情绪指标较前一周加速恶化，但非恐慌性抛售；港股关注度提升，南向资金流入加速。\par}
\paragraph{MS} PayPay评级从等权重上调至超配，目标价从24美元下调至23美元（反映寿险投资不确定性），但核心增长轨迹未变。预计未来三年收入CAGR 18.9\%，贡献利润CAGR 19.1\%；F3/27财年调整后EBITDA 1409亿日元，同比+26.8\%，高于公司指引中值。当前股价15.90美元，距目标价有约45\%上行空间。寿险投资（T\&D Financial）的不确定性已在估值折价中体现，核心业务（支付GMV扩张、变现率提升、金融服务交叉销售）依然强劲。
{\small\color{refgray}数据：目标价从24美元下调至23美元，评级从EW上调至OW\par}
{\small\color{refgray}数据：当前股价15.90美元，上行空间约45\%\par}
{\small\color{refgray}数据：F3/27财年调整后EBITDA预测1409亿日元（同比+26.8\%）\par}
{\small\color{refgray}数据：1季度调整后EBITDA预测322.5亿日元（公司指引305-325亿日元）\par}
{\small\color{refgray}边际变化：评级上调但目标价下调，反映市场对寿险投资的过度悲观已被定价，核心增长预期未变。\par}
\paragraph{NOM} BeOne 2Q26F收入预计16亿美元（同比+22\%），泽布替尼贡献11亿美元（同比+20\%），替雷利珠单抗2.13亿美元（同比+10\%）。毛利率88\%（+0.5pp），经营费用13亿美元（+19\%），经营利润1.52亿美元（+74\%）。但因1.29亿美元一次性补充税款，净利润仅5500万美元（同比-41\%）。全年收入目标63-65亿美元，毛利率88\%左右，GAAP经营利润7.5-8.5亿美元，下半年收入目标32亿美元。
{\small\color{refgray}数据：2Q26F收入16亿美元（同比+22\%）\par}
{\small\color{refgray}数据：泽布替尼销售11亿美元（同比+20\%，环比持平）\par}
{\small\color{refgray}数据：一次性税款1.29亿美元导致净利润同比-41\%至5500万美元\par}
{\small\color{refgray}数据：经营利润1.52亿美元（同比+74\%）\par}
{\small\color{refgray}边际变化：收入与经营利润高增长，但一次性税款造成净利润大幅下滑，市场需区分经营与财务节奏。\par}
\paragraph{NOM} 第一拖拉机股份从销量驱动转向产品组合与成本控制驱动。预计2026-2028年收入CAGR 6.90\%，净利润CAGR 12.48\%，利润增速持续快于收入。销量从2023年7.24万台降至2025年6.37万台，但均价从14.6万升至15.5万，毛利率从14.8\%升至15.1\%。大马力与智能化产品占比提升，海外收入2021-2025年CAGR约23\%，2025年海外收入占比11\%，预计2028年升至14\%。
{\small\color{refgray}数据：2026-2028年收入CAGR 6.90\%，净利润CAGR 12.48\%\par}
{\small\color{refgray}数据：销量从7.24万台（2023）降至6.37万台（2025）\par}
{\small\color{refgray}数据：均价从14.6万升至15.5万\par}
{\small\color{refgray}数据：毛利率从14.8\%（2023）升至15.1\%（2025），预计2028年达16.0\%\par}
{\small\color{refgray}边际变化：增长逻辑从销量驱动转向产品组合升级，利润增速持续快于收入，海外市场成为增量主要来源。\par}
\subsubsection*{关键数据}
\begin{itemize}
  \item MSASI加权读数从57\%降至37\%（7月15日至22日）
  \item 南向资金7月至今净流入约111亿美元
  \item PayPay当前股价15.90美元，距目标价23美元有约45\%上行空间
  \item BeOne 2Q26F经营利润+74\%但净利润-41\%，因1.29亿美元一次性税款
  \item 第一拖拉机股份2026-2028年净利润CAGR 12.48\%，快于收入CAGR 6.90\%
  \item 第一拖拉机股份海外收入2021-2025年CAGR约23\%
\end{itemize}
\begin{figure}[htbp]
\centering
\includegraphics[width=0.88\linewidth]{figures/fig_123.jpg}
\caption{Exhibit 1 - Exhibit 1: MS A-share Sentiment Indicator: MSASI weighted and MSASI weighted 1MMA Exhibit 2 : MSASI trajectory since January 1, 2019}
\end{figure}
\begin{figure}[htbp]
\centering
\includegraphics[width=0.88\linewidth]{figures/fig_128.jpg}
\caption{Exhibit 1 - Exhibit 1: PayPay scenario analysis \#\# Will the investment in T\&D Financial affect the company's growth trajectory? At least until the expected completion of the investment in October 2027, our view on PayPay's earnings outlook r}
\end{figure}
\begin{figure}[htbp]
\centering
\includegraphics[width=0.88\linewidth]{figures/fig_126.jpg}
\caption{Exhibit 4 - Exhibit 6: Equity futures turnover adjusted by moving 100D min-max (scaled to 0-100\% based on the percentage away from its 100-day high and low levels) vs. CSI 300 relative to 100D MA}
\end{figure}
\begin{figure}[htbp]
\centering
\includegraphics[width=0.88\linewidth]{figures/fig_131.jpg}
\caption{Exhibit 4 - Exhibit 7: T\&D Financial: MCEV has been growing steadily (JPY mn)}
\end{figure}
{\small\color{refgray}本节综合 4 篇研究；涉及 MS、NOM。}
\newpage
\section{辅助信号｜战略咨询}
本板块整合波士顿咨询近期研究，聚焦企业韧性对长期回报的贡献、成本管理优先级回升、欧洲太空探索的经济乘数效应、可再生能源扩张的核心驱动因素，以及AI在炼油行业的利润提升潜力。
\subsection{企业韧性：扰动时期的超额回报来源}
\textbf{扰动时期仅占研究周期的11\%，却解释了长期相对股东总回报（TSR）约30\%的变动，韧性是长期超额回报的关键来源。}
\begin{itemize}
  \item BCG基于近1800家美国公司25年数据发现，扰动季度内行业TSR分化程度（第25与75百分位差距）达30个百分点，是平稳期（16个百分点）的近两倍。
  \item 韧性拆解为三个维度：降低即时冲击（贡献63\%）、恢复速度（23\%）、恢复程度（15\%），其中抗压能力而非事后追赶最为关键。
  \item 15\%的企业具备“通用韧性”（在超80\%扰动季度跑赢行业），其25年年化TSR比行业高出5个百分点。
  \item 韧性并非成本，而是长期超额回报的重要来源，企业应重视扰动前的缓冲设计（如财务储备、供应链冗余、早期预警机制）。
\end{itemize}
\begin{figure}[htbp]
\centering
\includegraphics[width=0.88\linewidth]{figures/fig_148.jpg}
\caption{Exhibit 5 - Outperformance is driven primarily by withstanding the immediate impact. Of the three phases of resilience, not all are equally important. Reducing the immediate impact of a shock accounts for 63\% of relative TSR in crisis periods, while greater recovery speed}
\end{figure}
\subsection{成本管理优先级回升与执行短板}
\textbf{2025年成本管理取代增长扩张成为企业第一战略优先级，但企业平均仅实现48\%的成本节约目标，文化阻力是最大障碍。}
\begin{itemize}
  \item BCG对全球570余名高管调查显示，33\%将成本管理列为2025年第一优先级（较2024年上升8个百分点），70\%有足够中期可见度做出决策。
  \item 企业平均仅实现48\%的成本节约目标，未能达成目标的企业在股东总回报上比成功者落后9个百分点。
  \item 最大障碍是“员工与组织文化”层面的阻力，主动推动文化对齐的企业在成本削减效率上高出同行最多11个百分点。
  \item 供应链优化与产品组合简化是优先方向，成本管理本质是组织能力问题，而非单纯财务或运营技术。
\end{itemize}
\subsection{欧洲太空探索：500亿欧元投资的5倍经济乘数}
\textbf{若欧洲2025-2040年投入500亿欧元实施大型太空探索计划，累计GDP影响至少达2600亿欧元，乘数效应超5倍。}
\begin{itemize}
  \item 直接层面：新增约4000个全职岗位，产生250亿欧元直接GDP贡献；间接层面：采购支出产生约700亿欧元间接GDP影响；诱导层面：家庭消费贡献约550亿欧元。
  \item 催化效应（新太空市场）贡献额外1100亿欧元增量GDP，是拉开乘数差距的关键，主要来自发射、制造和在轨服务市场。
  \item 全球太空经济2022年规模约4600亿美元，赋能价值约3.1万亿美元；预计2040年分别达1万亿和7.9万亿美元。
  \item 欧洲目前仅占全球政府太空预算约15\%，与其约占全球25\%的GDP份额不匹配，不参与的成本可能不可忽视。
\end{itemize}
\begin{figure}[htbp]
\centering
\includegraphics[width=0.88\linewidth]{figures/fig_153.jpg}
\caption{Figure 2 - \#\# 3.1 Economic and societal benefits of space exploration Space exploration has the potential to create at least \textbackslash{}\textasciitilde{}€260 billion of cumulative GDP impact and an average of \textbackslash{}\textasciitilde{}90,000 FTEs in Europe between 2025 and 2040 (Figure 2). The investment in space explor}
\end{figure}
\subsection{可再生能源扩张：土地成本与装机基础是核心驱动}
\textbf{廉价土地和现有装机基础是中美两国风能和太阳能部署的首要驱动因素，其重要性超过资源禀赋或国家政策目标。}
\begin{itemize}
  \item BCG分析2014-2024年中美50个州和31个省份超10万个数据点，廉价土地解释了约20-21\%的装机增长，是最强预测变量。
  \item 美国土地成本最低的十个州风光装机平均增幅24\%，最高十个州仅9\%；新墨西哥州土地成本最低，实现40\%增长居全美之首。
  \item 现有装机基础解释了12-15\%的增长，体现经验曲线正向效应：早期部署降低后续供应链成本、施工经验和融资门槛。
  \item 美国需求增长解释了13\%的扩张，而中国因西部基地远离负荷中心，增长动力更多来自国家能源转型战略和跨区输电规划。
\end{itemize}
\begin{figure}[htbp]
\centering
\includegraphics[width=0.88\linewidth]{figures/fig_158.jpg}
\caption{Exhibit 2 - \#\# EXHIBIT 2 The Top Drivers Influencing Wind and Solar Scaling in the US and China Are Similar Relative importance of wind and solar adoption drivers across US states and Chinese provinces Note: Drivers are listed in descending order by combined share of scal}
\end{figure}
\subsection{AI在炼油行业：每桶0.5-1.2美元的利润提升}
\textbf{AI不再是可选项，而是炼油企业在产能过剩和利润波动中创造竞争优势的核心杠杆，中等水平玩家可获超50\%利润改善。}
\begin{itemize}
  \item 2026-2030年净炼油产能增长将超过需求增长，即使保守情景下供应仍超最乐观需求情景，利润承压。
  \item AI应用覆盖全价值链：规划与排程（每桶0.15-0.30美元）、运营（0.1-0.5美元）、物流与贸易（0.2-0.3美元）、维护与可靠性（0.05-0.1美元）。
  \item 东南亚炼油企业通过AI驱动的数字孪生进行短期排程优化，实现每桶0.15-0.30美元EBIT提升。
  \item 炼油企业需从周期性、孤立的优化转向实时、以利润为导向的决策闭环，覆盖规划、运营、物流、维护和安全等核心环节。
\end{itemize}
\newpage
\section{辅助信号｜智库/国际机构}
本板块整合经合组织两份最新报告，聚焦医疗数据互操作性与AI劳动力市场政策两大主题，揭示数据孤岛的经济代价、技术进展与治理短板，以及各国在AI技能培训、隐私保护与社会对话方面的政策分化。
\subsection{医疗数据互操作性：价值可观但治理瓶颈突出}
\textbf{经合组织报告估算，有效的数据互操作性每年可释放相当于卫生支出2.7\%至6.6\%的宏观价值，但技术标准进展未能解决治理、信任和运营层面的根本障碍。}
\begin{itemize}
  \item 加拿大、芬兰和英国等国的实证显示，数据互操作性每年可释放相当于卫生支出2.7\%至6.6\%的宏观环境价值，来源包括服务交付优化、错误减少、更及时的治疗和预防性干预。
  \item 当前结构性障碍显著：81.5\%的医生认为互操作性不足对患者安全构成潜在不确定性；患者因重复提供信息对医疗系统信任度下降15\%；诊断错误占医疗成本约17.5\%，其中相当比例与数据互操作不畅相关。
  \item 所有受访国家都在使用健康数据标准，41\%完全依赖国际标准，但标准采用不等于互操作性实现；受访者提及最多的障碍包括不一致的数据共享实践和去标识化做法（70.59\%）、区域间基础设施和技术碎片化（66.67\%），以及碎片化的治理和机构协调（60.78\%）。
  \item 报告识别出28种现行标准，涵盖词汇、内容结构、通信、隐私安全和标识符等领域，但技术层面的进展并未解决法律和运营层面的根本瓶颈，真正的挑战在于人、治理和领导力。
  \item 边际变化：经合组织首次系统量化互操作性的宏观价值区间，并强调治理短板而非技术标准是当前主要瓶颈，这为政策制定者提供了从投入产出角度推动系统性改革的依据。
\end{itemize}
\begin{figure}[htbp]
\centering
\includegraphics[width=0.88\linewidth]{figures/fig_140.jpg}
\caption{Figure 1 - \#\# What is interoperability in healthcare? 18. Interoperability is the ability of different health systems, devices, applications, or organisations to securely exchange, interpret, and use data in a coordinated and meaningful way, with minimal human interventi}
\end{figure}
\begin{figure}[htbp]
\centering
\includegraphics[width=0.88\linewidth]{figures/fig_143.jpg}
\caption{Figure 4 - 30. This analysis is guided by the Interoperability Tower (see Figure 4) which provides a frame of reference across local, regional, national, and cross-border levels of data exchange. With the complex nature of interoperability, and its common challenges invo}
\end{figure}
\subsection{AI劳动力市场政策：技能培训密集但覆盖偏窄，隐私与透明度仍处早期}
\textbf{经合组织对G7国家、欧盟及部分拉美地区的调查显示，多数国家未为AI单独设立劳动力市场政策，而是纳入现有框架；技能培训是政策最密集领域，但偏重技术能力而非AI素养，隐私、透明度和问责方面政策仍处于早期形成阶段。}
\begin{itemize}
  \item 所有被调查国家都提供再培训和技能提升项目，但大多数聚焦于技术技能而非通用AI素养，工人对AI如何影响自身工作、如何与AI系统协作的理解未被纳入主流培训目标。
  \item 针对因AI而失业的工人，各国主要依赖标准就业支持政策和社会保护，而非专门设计的AI转型支持；对于仍在职但面临AI冲击不确定性较高的工人，针对性措施仍然稀缺。
  \item 在隐私和数据保护方面，各国首先依赖现有法律框架，许多国家提供额外指南和工具（如工作坊、指导手册或自查工具）帮助雇主理解义务，但数据保护和隐私框架对工人及AI系统的具体影响并不总是被清晰界定或向雇主充分解释。
  \item 社会对话和集体谈判可帮助解决AI在劳动力市场使用中的实际挑战，但前提是各方具备相应的技术理解能力；目前多数国家在专门AI法规、AI素养培训和社会对话中的AI专业知识基础三个维度上仍有明显提升空间。
  \item 边际变化：报告提供了一个评估国家AI劳动力市场政策成熟度的三维框架（专门法规、素养培训、社会对话基础），并指出当前政策资源集中在“学会用AI工具”而非“理解AI如何改变工作逻辑”，后者对长期就业韧性更为关键。
\end{itemize}
\begin{figure}[htbp]
\centering
\includegraphics[width=0.88\linewidth]{figures/fig_145.jpg}
\caption{Figure 1 - At the same time, AI has the potential to help people with disabilities operate on the labour market. For instance, Touzet (2023[13]) interviewed over 70 stakeholders and identified 142 examples of AI-powered solutions to support people with disability, and em}
\end{figure}
\begin{figure}[htbp]
\centering
\includegraphics[width=0.88\linewidth]{figures/fig_146.jpg}
\caption{Figure 2 - \# 4 Non-discrimination in the labour market AI systems can help identify and address inequalities in the workplace by providing quantitative insights and supporting fairer decision making. For example, in the OECD AI surveys, 57\% of workers in manufacturing an}
\end{figure}
\section{结语}
后续需验证的关键节点包括：美联储7月FOMC会议声明及点阵图对利率路径的暗示；欧洲央行9月加息预期的兑现程度；Google等科技巨头AI资本开支的实际落地节奏与回报率；中国生物科技公司核心产品医保放量数据及海外临床进展；香港银行净息差持稳假设在利率下行压力下的可持续性；以及欧洲充电基础设施建设的政策响应与投资加速情况。
\newpage
\section{覆盖概览}
投行/券商 59 篇；战略咨询 5 篇；智库/国际机构 2 篇。\par
投行覆盖：NOM 15篇 / GS 14篇 / Bernstein 8篇 / JPM 6篇 / MS 6篇 / DB 5篇 / Citi 4篇 / BARC 1篇\par
\section{Disclaimer}
{\scriptsize\color{refgray}
This community is a paid learning and information-sharing community created for general educational purposes only. The membership fee is charged solely for access to community discussions, educational content, organization, and operational costs. It is not an advisory fee, management fee, performance fee, brokerage fee, or compensation for personalized financial, investment, legal, tax, or professional advice.\par
I participate and share information solely as an individual and for educational discussion among community members. I am not acting as a financial adviser, investment adviser, broker, dealer, portfolio manager, legal adviser, tax adviser, accountant, fiduciary, or any other licensed professional adviser. Nothing shared in this community constitutes, and should not be interpreted as, financial advice, investment advice, legal advice, tax advice, a recommendation, solicitation, offer, endorsement, instruction, or guarantee to buy, sell, hold, short, trade, or invest in any security, cryptocurrency, fund, derivative, strategy, or other financial product.\par
All content shared in this community, including messages, discussions, links, files, charts, examples, market commentary, watchlists, AI-generated or AI-polished summaries, and third-party materials, is general, impersonal, and educational in nature. It does not take into account any individual's financial situation, investment objectives, risk tolerance, investment horizon, liquidity needs, tax status, legal restrictions, or suitability requirements. No content should be relied upon as suitable for any specific person, account, portfolio, or financial decision.\par
Information shared in this community may be incomplete, inaccurate, outdated, speculative, AI-generated, AI-edited, or based on third-party internet sources that have not been independently verified. I make no representation or warranty as to the accuracy, completeness, timeliness, reliability, or suitability of any information shared.\par
Financial markets involve substantial risk, including the possible loss of principal. Past performance, historical data, hypothetical examples, simulated results, backtests, personal opinions, or educational examples do not guarantee future results. Each member is solely responsible for conducting their own research, exercising independent judgment, and making their own decisions. Before making any financial, investment, legal, or tax decision, members should consult qualified and licensed professionals.\par
Participation in this community, payment of any membership fee, or communication with me or other members does not create any adviser-client, fiduciary, professional, contractual, agency, partnership, employment, or other advisory relationship. By joining or participating in this community, you acknowledge and agree that you are solely responsible for your own decisions, actions, transactions, profits, losses, risks, and consequences, and that I shall not be liable for any direct, indirect, incidental, consequential, financial, legal, tax, or other loss or damage arising from or related to any content, discussion, information, or material shared in this community.\par
}
\newpage
\begin{center}
{\Large 更多详情报告kcdesk.com}\par\vspace{0.8cm}
\includegraphics[width=0.58\linewidth]{\detokenize{../../prompts/zsxq_img.jpg}}
\end{center}
\end{document}
